% TP POO Python — classes, attributs, méthodes, encapsulation, héritage (BTS CIEL)
% Compilation :
%   sujet élève : pdflatex TP1_POO_Python.tex
%   corrigé     : pdflatex -jobname=TP1_POO_Python_Corrige "\def\modecorrige{1}% TP POO Python — classes, attributs, méthodes, encapsulation, héritage (BTS CIEL)
% Compilation :
%   sujet élève : pdflatex TP1_POO_Python.tex
%   corrigé     : pdflatex -jobname=TP1_POO_Python_Corrige "\def\modecorrige{1}% TP POO Python — classes, attributs, méthodes, encapsulation, héritage (BTS CIEL)
% Compilation :
%   sujet élève : pdflatex TP1_POO_Python.tex
%   corrigé     : pdflatex -jobname=TP1_POO_Python_Corrige "\def\modecorrige{1}% TP POO Python — classes, attributs, méthodes, encapsulation, héritage (BTS CIEL)
% Compilation :
%   sujet élève : pdflatex TP1_POO_Python.tex
%   corrigé     : pdflatex -jobname=TP1_POO_Python_Corrige "\def\modecorrige{1}\input{TP1_POO_Python.tex}"
\documentclass[11pt,a4paper]{article}

\usepackage[T1]{fontenc}
\usepackage[utf8]{inputenc}
\usepackage{lmodern}
\usepackage[french]{babel}
\usepackage[a4paper,margin=1.8cm,bottom=2.2cm]{geometry}
\usepackage[table]{xcolor}
\usepackage{booktabs,tabularx,array}
\usepackage{enumitem}
\usepackage{needspace}
\usepackage{fancyhdr}
\usepackage{amssymb,amsmath}
\usepackage[most]{tcolorbox}
\usepackage{listings}
\usepackage[strings]{underscore}  % permet d'écrire __init__ sans échappement
\usepackage{hyperref}
\hypersetup{hidelinks}

\providecommand{\modecorrige}{0}
\newif\ifcorrige
\ifnum\modecorrige=1 \corrigetrue \else \corrigefalse \fi

% Palette alignée sur les autres supports du dépôt
\definecolor{ipred}{HTML}{C0392B}
\definecolor{cdcbar}{HTML}{A67C3D}
\definecolor{cdcbg}{HTML}{F4F1EA}
\definecolor{qbar}{HTML}{4A6B5D}
\definecolor{qbg}{HTML}{EDF1EF}
\definecolor{corr}{HTML}{1F5FA8}
\definecolor{corrbg}{HTML}{E8F0FA}
\definecolor{grisfoot}{HTML}{828282}
\definecolor{codebg}{HTML}{F7F7F4}
\definecolor{codekw}{HTML}{8E44AD}
\definecolor{codestr}{HTML}{1F7A4D}
\definecolor{codecom}{HTML}{8A8A8A}

% ---------------------------------------------------------------- code Python
\lstdefinestyle{py}{
  language=Python,
  basicstyle=\ttfamily\small,
  keywordstyle=\color{codekw}\bfseries,
  stringstyle=\color{codestr},
  commentstyle=\color{codecom}\itshape,
  showstringspaces=false,
  backgroundcolor=\color{codebg},
  frame=single,
  rulecolor=\color{black!25},
  framesep=5pt,
  xleftmargin=4pt,
  xrightmargin=4pt,
  columns=fullflexible,
  keepspaces=true,
  breaklines=true,
  extendedchars=true,
  literate={é}{{\'e}}1 {è}{{\`e}}1 {ê}{{\^e}}1 {à}{{\`a}}1 {ç}{{\c c}}1
           {ù}{{\`u}}1 {û}{{\^u}}1 {î}{{\^\i}}1 {ô}{{\^o}}1 {É}{{\'E}}1 {À}{{\`A}}1 {È}{{\`E}}1 {Ç}{{\c C}}1
           {«}{{\guillemotleft}}1 {»}{{\guillemotright}}1,
  morekeywords={self,True,False,None},
}
\lstset{style=py}
\newcommand{\py}[1]{\texttt{#1}}

% --------------------------------------------------------------- encadrés
\newtcolorbox{methode}[1][]{enhanced,breakable,colback=cdcbg,colframe=cdcbar,boxrule=0pt,
  leftrule=3pt,arc=0pt,outer arc=0pt,left=6pt,right=6pt,top=4pt,bottom=4pt,
  title={#1},coltitle=cdcbar,fonttitle=\bfseries,detach title,
  before upper={\tcbtitle\par\smallskip}}
\newtcolorbox{rappel}[1][]{enhanced,breakable,colback=qbg,colframe=qbar,boxrule=0pt,
  leftrule=3pt,arc=0pt,outer arc=0pt,left=6pt,right=6pt,top=4pt,bottom=4pt,
  title={#1},coltitle=qbar,fonttitle=\bfseries,attach title to upper={\ :\ }}
\newtcolorbox{piege}{enhanced,breakable,colback=red!4,colframe=ipred,boxrule=0pt,
  leftrule=3pt,arc=0pt,outer arc=0pt,left=6pt,right=6pt,top=4pt,bottom=4pt}

% --------------------------------------------------------------- exercices
\newcounter{exo}[section]
\renewcommand{\theexo}{\arabic{exo}}
\newtcolorbox{quest}{enhanced,breakable,colback=qbg,colframe=qbar,boxrule=0pt,
  leftrule=3pt,arc=0pt,outer arc=0pt,left=6pt,right=6pt,top=4pt,bottom=4pt}
\newcommand{\exo}[2][]{% [sous-titre] {énoncé}
  \Needspace{5cm}\stepcounter{exo}%
  \begin{quest}\textbf{Exercice \theexo\ifx\relax#1\relax\else\ — #1\fi.}
  \ifx\relax#2\relax\else\ #2\fi\end{quest}}

% Réponse courte : pointillés dans le sujet, texte bleu dans le corrigé
\newcommand{\rep}[1]{%
  \ifcorrige
    \par\smallskip{\color{corr}\small#1}\par\smallskip
  \else
    \par\nobreak\vspace{1mm}\noindent\color{black!55}\dotfill\par
    \noindent\color{black!55}\dotfill\par\color{black}\vspace{1mm}
  \fi}

% Zone de code vierge (sujet uniquement) de hauteur donnée
\newcommand{\zonecode}[1][4cm]{%
  \ifcorrige\else
    \par\smallskip\noindent
    \begin{tcolorbox}[colback=white,colframe=black!45,boxrule=0.5pt,arc=1pt,height=#1]\end{tcolorbox}%
  \fi}

% Bloc visible seulement dans le corrigé
\newenvironment{solution}
  {\ifcorrige
     \begin{tcolorbox}[enhanced,breakable,colback=corrbg,colframe=corr,boxrule=0.6pt,arc=1pt,
       left=5pt,right=5pt,top=4pt,bottom=4pt,
       title={Corrigé},coltitle=corr,fonttitle=\bfseries\small,detach title,
       before upper={\tcbtitle\par\smallskip},fontupper=\small]%
   \else
     \setbox0\vbox\bgroup
   \fi}
  {\ifcorrige\end{tcolorbox}\else\egroup\fi}

\pagestyle{fancy}
\fancyhf{}
\renewcommand{\headrulewidth}{0pt}
\fancyfoot[L]{\small\itshape\color{grisfoot}TP POO Python\ifcorrige\ — CORRIGÉ\fi}
\fancyfoot[C]{\small\color{grisfoot}BTS CIEL}
\fancyfoot[R]{\small\color{grisfoot}\thepage}

\setlist{itemsep=2pt,topsep=3pt}
\setlength{\parindent}{0pt}
\setlength{\parskip}{4pt}
\renewcommand{\arraystretch}{1.3}

\begin{document}

\begin{center}
{\LARGE\bfseries Programmation orientée objet en Python\par}
\smallskip
{\large Classes, attributs, méthodes --- Encapsulation et héritage\par}
\medskip
BTS CIEL --- Travaux pratiques
\ifcorrige\par\textcolor{corr}{\textbf{Version enseignant --- corrigés intégrés}}\fi
\end{center}

\medskip
\noindent\begin{tabularx}{\textwidth}{@{}>{\bfseries}l X@{}}
Nom, prénom : & \ifcorrige\textcolor{corr}{corrigé}\else\dotfill\fi \\
Classe : & \ifcorrige\textcolor{corr}{---}\else\dotfill\fi \\
Date : & \ifcorrige\textcolor{corr}{---}\else\dotfill\fi \\
\end{tabularx}

\begin{methode}[Objectifs du TP]
\begin{itemize}
  \item Lire une classe Python et identifier son constructeur, ses attributs et ses méthodes.
  \item Instancier un objet et appeler ses méthodes.
  \item Compléter une classe existante en ajoutant une méthode.
  \item Protéger les données d'un objet par l'encapsulation (attributs privés, accesseurs, mutateurs).
  \item Écrire un mutateur qui valide la donnée avant de l'enregistrer.
  \item Construire une classe fille par héritage, redéfinir une méthode et appeler le constructeur parent.
\end{itemize}
\textbf{Consignes} : chaque question se traite d'abord sur le papier, puis se vérifie dans
l'interpréteur Python (fichier \py{tp\_poo.py}). Le code rendu doit s'exécuter sans erreur.
\end{methode}

\begin{rappel}[Vocabulaire de référence]
Une \textbf{classe} est un modèle. Un \textbf{objet} (ou \textbf{instance}) est une réalisation
concrète de ce modèle. Un \textbf{attribut} est une donnée portée par l'objet. Une
\textbf{méthode} est une fonction définie dans la classe et appelée sur un objet. Le paramètre
\py{self} désigne l'objet sur lequel la méthode est appelée.
\end{rappel}

% ============================================================================
\section*{Partie A --- Classes, attributs et méthodes}
\addcontentsline{toc}{section}{Partie A}
\setcounter{exo}{0}

% ---------------------------------------------------------------- A1
\exo[La classe \py{Voiture}]{Lire le code suivant, puis répondre aux questions.}

\begin{lstlisting}
class Voiture:
    def __init__(self, marque, couleur):
        self.marque = marque
        self.couleur = couleur

    def afficher_details(self):
        return f"Voiture {self.marque} de couleur {self.couleur}"
\end{lstlisting}

\textbf{Questions}
\begin{enumerate}
  \item Que fait la méthode \py{__init__} ?
    \rep{C'est le \textbf{constructeur}. Il est appelé automatiquement à la création de l'objet
    et initialise les attributs d'instance \py{self.marque} et \py{self.couleur} à partir des
    valeurs passées en paramètres.}
  \item Combien d'attributs possède la classe \py{Voiture} ?
    \rep{Deux attributs d'instance : \py{marque} et \py{couleur}.}
  \item Que retourne la méthode \py{afficher_details} ?
    \rep{Une chaîne de caractères (\py{str}) construite par une f-string, par exemple
    \py{"Voiture Toyota de couleur Bleu"}. Elle \emph{retourne} la chaîne, elle ne l'affiche pas.}
\end{enumerate}

\textbf{À compléter} : créer une instance de \py{Voiture} de marque \py{"Toyota"} et de couleur
\py{"Bleu"}, appeler \py{afficher_details} et afficher le résultat.
\zonecode[3cm]
\begin{solution}
\begin{lstlisting}
ma_voiture = Voiture("Toyota", "Bleu")
print(ma_voiture.afficher_details())
# Voiture Toyota de couleur Bleu
\end{lstlisting}
\py{print} est nécessaire : sans lui, la chaîne est retournée puis perdue.
\end{solution}

% ---------------------------------------------------------------- A2
\exo[Compléter la classe \py{Chien}]{}

\begin{lstlisting}
class Chien:
    def __init__(self, nom, age):
        self.nom = nom
        self.age = age

    # Méthode aboyer qui retourne "Wouf !"
\end{lstlisting}

\textbf{Questions}
\begin{enumerate}
  \item Quels sont les attributs de la classe \py{Chien} ?
    \rep{\py{nom} et \py{age}, deux attributs d'instance.}
  \item Comment appelle-t-on la méthode \py{__init__} dans le contexte de la POO ?
    \rep{Le \textbf{constructeur} (ou initialiseur). C'est une \emph{méthode spéciale}, reconnaissable
    à son double souligné ; on ne l'appelle jamais directement, elle est déclenchée par
    \py{Chien("Rex", 3)}.}
  \item Ajouter une méthode \py{aboyer} qui retourne \py{"Wouf !"}.
  \item Créer un objet \py{Chien} nommé \py{"Rex"}, âgé de 3 ans, et le faire aboyer.
\end{enumerate}
\zonecode[4.5cm]
\begin{solution}
\begin{lstlisting}
class Chien:
    def __init__(self, nom, age):
        self.nom = nom
        self.age = age

    def aboyer(self):
        return "Wouf !"

rex = Chien("Rex", 3)
print(rex.aboyer())   # Wouf !
\end{lstlisting}
Le paramètre \py{self} est obligatoire dans la définition, même si \py{aboyer} n'utilise aucun
attribut : il reçoit l'objet sur lequel la méthode est appelée.
\end{solution}

% ---------------------------------------------------------------- A3
\exo[La classe \py{Livre}]{}

\begin{lstlisting}
class Livre:
    def __init__(self, titre, auteur):
        self.titre = titre
        self.auteur = auteur

    def description(self):
        return f"{self.titre} par {self.auteur}"
\end{lstlisting}

\textbf{Questions}
\begin{enumerate}
  \item Que fait la méthode \py{description} ?
    \rep{Elle retourne une chaîne qui associe le titre et l'auteur du livre, sous la forme
    \py{"titre par auteur"}.}
  \item Combien de paramètres la méthode \py{__init__} attend-elle, en incluant \py{self} ?
    \rep{Trois : \py{self}, \py{titre} et \py{auteur}. À l'appel, seuls deux arguments sont
    fournis : \py{self} est transmis automatiquement par Python.}
  \item Créer un objet \py{Livre} de titre \py{"Le Petit Prince"} et d'auteur
    \py{"Antoine de Saint-Exupéry"}, puis afficher sa description.
\end{enumerate}
\zonecode[3cm]
\begin{solution}
\begin{lstlisting}
livre = Livre("Le Petit Prince", "Antoine de Saint-Exupery")
print(livre.description())
# Le Petit Prince par Antoine de Saint-Exupery
\end{lstlisting}
\end{solution}

% ---------------------------------------------------------------- A4
\exo[La classe \py{Rectangle}]{}

\begin{lstlisting}
class Rectangle:
    def __init__(self, longueur, largeur):
        self.longueur = longueur
        self.largeur = largeur

    # Méthode calculer_aire qui retourne l'aire du rectangle
\end{lstlisting}

\textbf{Questions}
\begin{enumerate}
  \item Quels sont les attributs de la classe \py{Rectangle} ?
    \rep{\py{longueur} et \py{largeur}.}
  \item Quelle formule mathématique faut-il utiliser pour calculer l'aire ?
    \rep{$\text{aire} = \text{longueur} \times \text{largeur}$.}
  \item Ajouter une méthode \py{calculer_aire} qui retourne l'aire du rectangle.
  \item Créer un rectangle de longueur 5 et de largeur 3, puis afficher son aire.
\end{enumerate}
\zonecode[4.5cm]
\begin{solution}
\begin{lstlisting}
class Rectangle:
    def __init__(self, longueur, largeur):
        self.longueur = longueur
        self.largeur = largeur

    def calculer_aire(self):
        return self.longueur * self.largeur

r = Rectangle(5, 3)
print(r.calculer_aire())   # 15
\end{lstlisting}
Une méthode peut aussi être ajoutée pour le périmètre :\\
\py{return 2 * (self.longueur + self.largeur)}.
\end{solution}

% ---------------------------------------------------------------- A5
\exo[Création complète d'une classe \py{Etudiant}]{On souhaite modéliser un étudiant par son
nom et sa note.}

\begin{enumerate}
  \item Créer une classe \py{Etudiant} possédant un constructeur \py{__init__} prenant en
    paramètres \py{nom} et \py{note}, et deux attributs d'instance \py{self.nom} et \py{self.note}.
  \item Ajouter une méthode \py{a_reussi} qui retourne \py{True} si la note est supérieure ou
    égale à 10, et \py{False} sinon.
  \item Tester la classe en créant un étudiant avec une note de 8.
  \item Ajouter une méthode \py{augmenter_note} qui augmente la note de 2 points.
  \item Créer un étudiant nommé \py{"Alice"} avec une note de 15 et afficher si elle a réussi,
    grâce à la méthode \py{a_reussi}.
\end{enumerate}
\zonecode[7cm]
\begin{solution}
\begin{lstlisting}
class Etudiant:
    def __init__(self, nom, note):
        self.nom = nom
        self.note = note

    def a_reussi(self):
        return self.note >= 10

    def augmenter_note(self):
        self.note = self.note + 2

# 3. test avec une note de 8
e = Etudiant("Bob", 8)
print(e.a_reussi())          # False

# 5. Alice avec une note de 15
alice = Etudiant("Alice", 15)
print(alice.a_reussi())      # True
\end{lstlisting}
\textbf{Points d'attention.}
\py{return self.note >= 10} suffit : la comparaison produit déjà un booléen, un
\py{if ... return True else return False} est inutile.
\py{augmenter_note} \emph{modifie} l'attribut et ne retourne rien : elle agit sur l'état de l'objet.
\end{solution}

% ============================================================================
\clearpage
\section*{Partie B --- Encapsulation et héritage}
\addcontentsline{toc}{section}{Partie B}
\setcounter{exo}{0}

\begin{rappel}[Encapsulation]
Préfixer un attribut par deux soulignés (\py{self.__nom}) déclenche le \emph{name mangling} :
Python renomme l'attribut en \py{_Classe__nom}, ce qui le rend inaccessible par son nom
d'origine depuis l'extérieur. On y accède alors par un \textbf{accesseur} (\emph{getter},
\py{get_nom}) et on le modifie par un \textbf{mutateur} (\emph{setter}, \py{set_nom}), ce qui
permet de contrôler les valeurs acceptées.
\end{rappel}

\begin{rappel}[Héritage]
\py{class Chat(Animal)} signifie que \py{Chat} \textbf{hérite} de \py{Animal} : la classe fille
reçoit les attributs et méthodes de la classe mère. Elle peut en ajouter, ou
\textbf{redéfinir} une méthode existante (\emph{override}). \py{super().__init__(...)} appelle le
constructeur de la classe mère.
\end{rappel}

% ---------------------------------------------------------------- B1
\exo[Encapsulation simple avec attribut privé]{}

\begin{lstlisting}
class Personne:
    def __init__(self, nom, age):
        self.__nom = nom
        self.__age = age

    def get_nom(self):
        return self.__nom

    # Méthode get_age()
\end{lstlisting}

\textbf{Questions}
\begin{enumerate}
  \item Pourquoi les attributs \py{__nom} et \py{__age} sont-ils privés ?
    \rep{Pour empêcher leur modification directe depuis l'extérieur de la classe. L'accès passe
    par des méthodes, ce qui permet de contrôler la lecture et l'écriture et de modifier
    l'implémentation interne sans casser le code qui utilise la classe.}
  \item Comment accède-t-on à l'attribut \py{__nom} ?
    \rep{Par l'accesseur : \py{p.get_nom()}. L'écriture \py{p.__nom} provoque une
    \py{AttributeError}.}
  \item Ajouter une méthode \py{get_age} pour accéder à l'attribut \py{__age}.
  \item Créer une instance de \py{Personne} de nom \py{"Luc"} et d'âge 25, puis afficher son nom.
\end{enumerate}
\zonecode[4.5cm]
\begin{solution}
\begin{lstlisting}
class Personne:
    def __init__(self, nom, age):
        self.__nom = nom
        self.__age = age

    def get_nom(self):
        return self.__nom

    def get_age(self):
        return self.__age

p = Personne("Luc", 25)
print(p.get_nom())    # Luc
print(p.get_age())    # 25
\end{lstlisting}
\end{solution}

% ---------------------------------------------------------------- B2
\exo[Ajouter un mutateur]{}

\begin{lstlisting}
class CompteBancaire:
    def __init__(self, titulaire, solde):
        self.__titulaire = titulaire
        self.__solde = solde

    def get_solde(self):
        return self.__solde
\end{lstlisting}

\textbf{Questions}
\begin{enumerate}
  \item Que signifie l'encapsulation dans cette classe ?
    \rep{Le titulaire et le solde sont privés : aucun code extérieur ne peut écrire
    \py{compte.__solde = 1000000}. Le solde ne peut être consulté que par \py{get_solde} et
    modifié que par les méthodes que la classe met à disposition.}
  \item Que fait la méthode \py{get_solde} ?
    \rep{Elle retourne la valeur de l'attribut privé \py{__solde} : c'est un accesseur en
    lecture seule.}
  \item Créer un compte bancaire pour \py{"Emma"} avec un solde de 500 et afficher le solde.
  \item Ajouter une méthode \py{set_solde} pour modifier le solde.
\end{enumerate}
\zonecode[5cm]
\begin{solution}
\begin{lstlisting}
class CompteBancaire:
    def __init__(self, titulaire, solde):
        self.__titulaire = titulaire
        self.__solde = solde

    def get_solde(self):
        return self.__solde

    def set_solde(self, nouveau_solde):
        self.__solde = nouveau_solde

compte = CompteBancaire("Emma", 500)
print(compte.get_solde())   # 500
compte.set_solde(750)
print(compte.get_solde())   # 750
\end{lstlisting}
\textbf{Prolongement.} Un mutateur brut n'apporte pas de sécurité. Une version plus réaliste
refuse un solde négatif, ou expose plutôt des opérations métier :
\py{deposer(montant)} et \py{retirer(montant)} avec vérification de la provision.
\end{solution}

% ---------------------------------------------------------------- B3
\exo[Encapsulation avec validation]{}

\begin{lstlisting}
class Etudiant:
    def __init__(self, nom, note):
        self.__nom = nom
        self.__note = note

    def get_note(self):
        return self.__note

    # Méthode set_note qui accepte uniquement des notes entre 0 et 20
\end{lstlisting}

\textbf{Questions}
\begin{enumerate}
  \item Pourquoi utiliser un mutateur pour la note ?
    \rep{Pour garantir un \emph{invariant} : une note reste comprise entre 0 et 20. Le mutateur
    est le seul point d'entrée en écriture, il peut donc vérifier la valeur avant de
    l'enregistrer et refuser les valeurs incohérentes.}
  \item Que se passe-t-il si on essaie d'accéder directement à \py{__note} ?
    \rep{Python lève une \py{AttributeError} : à cause du \emph{name mangling}, l'attribut
    s'appelle en réalité \py{_Etudiant__note} et le nom \py{__note} n'existe pas depuis
    l'extérieur.}
  \item Ajouter une méthode \py{set_note} qui n'accepte que des notes entre 0 et 20.
  \item Créer un étudiant nommé \py{"Paul"} avec une note de 15, puis essayer de changer la note
    à 25.
\end{enumerate}
\zonecode[5.5cm]
\begin{solution}
\begin{lstlisting}
class Etudiant:
    def __init__(self, nom, note):
        self.__nom = nom
        self.__note = note

    def get_note(self):
        return self.__note

    def set_note(self, nouvelle_note):
        if 0 <= nouvelle_note <= 20:
            self.__note = nouvelle_note
        else:
            print("Note invalide : elle doit etre comprise entre 0 et 20.")

paul = Etudiant("Paul", 15)
print(paul.get_note())   # 15
paul.set_note(25)        # Note invalide : ...
print(paul.get_note())   # 15 : la note n'a pas change
\end{lstlisting}
\textbf{Variante professionnelle.} Au lieu d'un \py{print}, lever une exception rend l'erreur
impossible à ignorer : \py{raise ValueError("note hors intervalle")}.
\end{solution}

% ---------------------------------------------------------------- B4
\exo[Héritage simple]{}

\begin{lstlisting}
class Animal:
    def __init__(self, nom):
        self.__nom = nom

    def get_nom(self):
        return self.__nom

    def faire_son(self):
        return "Son generique"

class Chat(Animal):
    # Méthode faire_son pour retourner "Miaou"
\end{lstlisting}

\textbf{Questions}
\begin{enumerate}
  \item Quelle classe est la classe parent ?
    \rep{\py{Animal}. \py{Chat} est la classe fille (ou classe dérivée).}
  \item Que fait l'héritage dans ce cas ?
    \rep{\py{Chat} récupère sans les réécrire le constructeur \py{__init__}, l'attribut privé
    \py{__nom} et la méthode \py{get_nom} d'\py{Animal}. Elle peut en plus redéfinir
    \py{faire_son} pour lui donner son comportement propre.}
  \item Redéfinir la méthode \py{faire_son} pour qu'elle retourne \py{"Miaou"}.
  \item Créer un \py{Chat} nommé \py{"Minou"} et afficher son son.
\end{enumerate}
\zonecode[4.5cm]
\begin{solution}
\begin{lstlisting}
class Chat(Animal):
    def faire_son(self):
        return "Miaou"

minou = Chat("Minou")
print(minou.get_nom())     # Minou  (methode heritee)
print(minou.faire_son())   # Miaou  (methode redefinie)
\end{lstlisting}
\py{Chat} ne définit pas de \py{__init__} : celui d'\py{Animal} est utilisé tel quel.
Appeler \py{faire_son()} sur un \py{Chat} exécute la version de \py{Chat} : c'est le
\textbf{polymorphisme}.
\end{solution}

% ---------------------------------------------------------------- B5
\exo[Héritage avec nouvel attribut]{}

\begin{lstlisting}
class Vehicule:
    def __init__(self, marque):
        self.__marque = marque

    def get_marque(self):
        return self.__marque

class Voiture(Vehicule):
    def __init__(self, marque, couleur):
        super().__init__(marque)
        self.__couleur = couleur

    # À compléter : ajoutez une méthode get_couleur
\end{lstlisting}

\textbf{Questions}
\begin{enumerate}
  \item Que fait \py{super().__init__(marque)} ?
    \rep{Il appelle le constructeur de la classe mère \py{Vehicule} pour initialiser la partie
    héritée de l'objet, ici l'attribut privé \py{__marque}. Sans cet appel, \py{get_marque}
    échouerait car l'attribut n'aurait jamais été créé.}
  \item Combien d'attributs privés possède la classe \py{Voiture} ?
    \rep{Deux au total dans l'objet : \py{_Vehicule__marque}, hérité et initialisé par
    \py{super()}, et \py{_Voiture__couleur}, propre à \py{Voiture}. La classe \py{Voiture} n'en
    déclare qu'un seul elle-même.}
  \item À compléter : ajouter une méthode \py{get_couleur}.
\end{enumerate}
\zonecode[4cm]
\begin{solution}
\begin{lstlisting}
class Voiture(Vehicule):
    def __init__(self, marque, couleur):
        super().__init__(marque)
        self.__couleur = couleur

    def get_couleur(self):
        return self.__couleur

v = Voiture("Renault", "Bleu")
print(v.get_marque())    # Renault
print(v.get_couleur())   # Bleu
\end{lstlisting}
\begin{piege}
\py{__marque} est privé \emph{à \py{Vehicule}} : à cause du \emph{name mangling},
une méthode de \py{Voiture} ne peut pas écrire \py{self.__marque}. Elle doit passer par
\py{self.get_marque()}. C'est exactement le rôle de l'accesseur hérité.
\end{piege}
\end{solution}

% ---------------------------------------------------------------- B6
\exo[Encapsulation et mutateur avec validation]{}

\begin{lstlisting}
class Livre:
    def __init__(self, titre, prix):
        self.__titre = titre
        self.__prix = prix

    def get_prix(self):
        return self.__prix

    # À compléter : ajoutez une méthode set_prix qui refuse les prix négatifs
\end{lstlisting}

\textbf{Questions}
\begin{enumerate}
  \item Pourquoi le prix est-il un attribut privé ?
    \rep{Parce qu'un prix doit respecter une règle métier : il ne peut pas être négatif. Le
    rendre privé interdit l'écriture directe et oblige à passer par une méthode qui applique
    cette règle.}
  \item Que fait un mutateur avec validation ?
    \rep{Il vérifie la valeur reçue avant de l'affecter à l'attribut. Si la valeur est valide,
    il met l'attribut à jour ; sinon il refuse la modification (message d'erreur ou exception)
    et l'objet conserve son état précédent, qui reste cohérent.}
  \item Créer un \py{Livre} intitulé \py{"Python 101"} avec un prix de 30, puis essayer de
    définir un prix de $-10$.
\end{enumerate}
\zonecode[5.5cm]
\begin{solution}
\begin{lstlisting}
class Livre:
    def __init__(self, titre, prix):
        self.__titre = titre
        self.__prix = prix

    def get_prix(self):
        return self.__prix

    def set_prix(self, nouveau_prix):
        if nouveau_prix >= 0:
            self.__prix = nouveau_prix
        else:
            print("Prix invalide : un prix ne peut pas etre negatif.")

livre = Livre("Python 101", 30)
print(livre.get_prix())   # 30
livre.set_prix(-10)       # Prix invalide : ...
print(livre.get_prix())   # 30 : le prix est inchange
\end{lstlisting}
\textbf{Remarque.} Le constructeur devrait appliquer la même validation, par exemple en
appelant \py{self.set_prix(prix)} au lieu d'affecter directement : sinon
\py{Livre("X", -10)} crée un objet invalide dès le départ.
\end{solution}

% ============================================================================
\clearpage
\section*{Synthèse}

\begin{center}
\begin{tabularx}{\textwidth}{@{}>{\bfseries}p{4.3cm} X@{}}
\toprule
Notion & Ce qu'il faut retenir \\
\midrule
Classe & Modèle décrivant des données et des comportements : \py{class Voiture:} \\
Objet / instance & Réalisation concrète de la classe : \py{v = Voiture("Toyota", "Bleu")} \\
Constructeur & \py{__init__}, appelé automatiquement, initialise les attributs \\
\py{self} & Référence à l'objet courant ; premier paramètre de toute méthode d'instance \\
Attribut d'instance & Donnée propre à l'objet : \py{self.marque} \\
Méthode & Fonction définie dans la classe, appelée par \py{objet.methode()} \\
Attribut privé & \py{self.__prix} : inaccessible depuis l'extérieur (\emph{name mangling}) \\
Accesseur (\emph{getter}) & \py{get_prix()} : lit un attribut privé \\
Mutateur (\emph{setter}) & \py{set_prix(v)} : écrit un attribut privé après validation \\
Héritage & \py{class Voiture(Vehicule)} : la fille reçoit les membres de la mère \\
\py{super()} & Accède à la classe mère, notamment \py{super().__init__(...)} \\
Redéfinition & La fille réécrit une méthode de la mère (\py{faire_son}) \\
\bottomrule
\end{tabularx}
\end{center}

\begin{piege}
\textbf{Erreurs les plus fréquentes.}
oublier \py{self} dans la définition d'une méthode ;
confondre \py{return} (retourne une valeur) et \py{print} (affiche) ;
écrire \py{objet.__attribut} depuis l'extérieur ;
oublier \py{super().__init__(...)} dans le constructeur d'une classe fille ;
appeler une méthode sans les parenthèses (\py{v.afficher_details} retourne la méthode,
pas son résultat).
\end{piege}

\end{document}
"
\documentclass[11pt,a4paper]{article}

\usepackage[T1]{fontenc}
\usepackage[utf8]{inputenc}
\usepackage{lmodern}
\usepackage[french]{babel}
\usepackage[a4paper,margin=1.8cm,bottom=2.2cm]{geometry}
\usepackage[table]{xcolor}
\usepackage{booktabs,tabularx,array}
\usepackage{enumitem}
\usepackage{needspace}
\usepackage{fancyhdr}
\usepackage{amssymb,amsmath}
\usepackage[most]{tcolorbox}
\usepackage{listings}
\usepackage[strings]{underscore}  % permet d'écrire __init__ sans échappement
\usepackage{hyperref}
\hypersetup{hidelinks}

\providecommand{\modecorrige}{0}
\newif\ifcorrige
\ifnum\modecorrige=1 \corrigetrue \else \corrigefalse \fi

% Palette alignée sur les autres supports du dépôt
\definecolor{ipred}{HTML}{C0392B}
\definecolor{cdcbar}{HTML}{A67C3D}
\definecolor{cdcbg}{HTML}{F4F1EA}
\definecolor{qbar}{HTML}{4A6B5D}
\definecolor{qbg}{HTML}{EDF1EF}
\definecolor{corr}{HTML}{1F5FA8}
\definecolor{corrbg}{HTML}{E8F0FA}
\definecolor{grisfoot}{HTML}{828282}
\definecolor{codebg}{HTML}{F7F7F4}
\definecolor{codekw}{HTML}{8E44AD}
\definecolor{codestr}{HTML}{1F7A4D}
\definecolor{codecom}{HTML}{8A8A8A}

% ---------------------------------------------------------------- code Python
\lstdefinestyle{py}{
  language=Python,
  basicstyle=\ttfamily\small,
  keywordstyle=\color{codekw}\bfseries,
  stringstyle=\color{codestr},
  commentstyle=\color{codecom}\itshape,
  showstringspaces=false,
  backgroundcolor=\color{codebg},
  frame=single,
  rulecolor=\color{black!25},
  framesep=5pt,
  xleftmargin=4pt,
  xrightmargin=4pt,
  columns=fullflexible,
  keepspaces=true,
  breaklines=true,
  extendedchars=true,
  literate={é}{{\'e}}1 {è}{{\`e}}1 {ê}{{\^e}}1 {à}{{\`a}}1 {ç}{{\c c}}1
           {ù}{{\`u}}1 {û}{{\^u}}1 {î}{{\^\i}}1 {ô}{{\^o}}1 {É}{{\'E}}1 {À}{{\`A}}1 {È}{{\`E}}1 {Ç}{{\c C}}1
           {«}{{\guillemotleft}}1 {»}{{\guillemotright}}1,
  morekeywords={self,True,False,None},
}
\lstset{style=py}
\newcommand{\py}[1]{\texttt{#1}}

% --------------------------------------------------------------- encadrés
\newtcolorbox{methode}[1][]{enhanced,breakable,colback=cdcbg,colframe=cdcbar,boxrule=0pt,
  leftrule=3pt,arc=0pt,outer arc=0pt,left=6pt,right=6pt,top=4pt,bottom=4pt,
  title={#1},coltitle=cdcbar,fonttitle=\bfseries,detach title,
  before upper={\tcbtitle\par\smallskip}}
\newtcolorbox{rappel}[1][]{enhanced,breakable,colback=qbg,colframe=qbar,boxrule=0pt,
  leftrule=3pt,arc=0pt,outer arc=0pt,left=6pt,right=6pt,top=4pt,bottom=4pt,
  title={#1},coltitle=qbar,fonttitle=\bfseries,attach title to upper={\ :\ }}
\newtcolorbox{piege}{enhanced,breakable,colback=red!4,colframe=ipred,boxrule=0pt,
  leftrule=3pt,arc=0pt,outer arc=0pt,left=6pt,right=6pt,top=4pt,bottom=4pt}

% --------------------------------------------------------------- exercices
\newcounter{exo}[section]
\renewcommand{\theexo}{\arabic{exo}}
\newtcolorbox{quest}{enhanced,breakable,colback=qbg,colframe=qbar,boxrule=0pt,
  leftrule=3pt,arc=0pt,outer arc=0pt,left=6pt,right=6pt,top=4pt,bottom=4pt}
\newcommand{\exo}[2][]{% [sous-titre] {énoncé}
  \Needspace{5cm}\stepcounter{exo}%
  \begin{quest}\textbf{Exercice \theexo\ifx\relax#1\relax\else\ — #1\fi.}
  \ifx\relax#2\relax\else\ #2\fi\end{quest}}

% Réponse courte : pointillés dans le sujet, texte bleu dans le corrigé
\newcommand{\rep}[1]{%
  \ifcorrige
    \par\smallskip{\color{corr}\small#1}\par\smallskip
  \else
    \par\nobreak\vspace{1mm}\noindent\color{black!55}\dotfill\par
    \noindent\color{black!55}\dotfill\par\color{black}\vspace{1mm}
  \fi}

% Zone de code vierge (sujet uniquement) de hauteur donnée
\newcommand{\zonecode}[1][4cm]{%
  \ifcorrige\else
    \par\smallskip\noindent
    \begin{tcolorbox}[colback=white,colframe=black!45,boxrule=0.5pt,arc=1pt,height=#1]\end{tcolorbox}%
  \fi}

% Bloc visible seulement dans le corrigé
\newenvironment{solution}
  {\ifcorrige
     \begin{tcolorbox}[enhanced,breakable,colback=corrbg,colframe=corr,boxrule=0.6pt,arc=1pt,
       left=5pt,right=5pt,top=4pt,bottom=4pt,
       title={Corrigé},coltitle=corr,fonttitle=\bfseries\small,detach title,
       before upper={\tcbtitle\par\smallskip},fontupper=\small]%
   \else
     \setbox0\vbox\bgroup
   \fi}
  {\ifcorrige\end{tcolorbox}\else\egroup\fi}

\pagestyle{fancy}
\fancyhf{}
\renewcommand{\headrulewidth}{0pt}
\fancyfoot[L]{\small\itshape\color{grisfoot}TP POO Python\ifcorrige\ — CORRIGÉ\fi}
\fancyfoot[C]{\small\color{grisfoot}BTS CIEL}
\fancyfoot[R]{\small\color{grisfoot}\thepage}

\setlist{itemsep=2pt,topsep=3pt}
\setlength{\parindent}{0pt}
\setlength{\parskip}{4pt}
\renewcommand{\arraystretch}{1.3}

\begin{document}

\begin{center}
{\LARGE\bfseries Programmation orientée objet en Python\par}
\smallskip
{\large Classes, attributs, méthodes --- Encapsulation et héritage\par}
\medskip
BTS CIEL --- Travaux pratiques
\ifcorrige\par\textcolor{corr}{\textbf{Version enseignant --- corrigés intégrés}}\fi
\end{center}

\medskip
\noindent\begin{tabularx}{\textwidth}{@{}>{\bfseries}l X@{}}
Nom, prénom : & \ifcorrige\textcolor{corr}{corrigé}\else\dotfill\fi \\
Classe : & \ifcorrige\textcolor{corr}{---}\else\dotfill\fi \\
Date : & \ifcorrige\textcolor{corr}{---}\else\dotfill\fi \\
\end{tabularx}

\begin{methode}[Objectifs du TP]
\begin{itemize}
  \item Lire une classe Python et identifier son constructeur, ses attributs et ses méthodes.
  \item Instancier un objet et appeler ses méthodes.
  \item Compléter une classe existante en ajoutant une méthode.
  \item Protéger les données d'un objet par l'encapsulation (attributs privés, accesseurs, mutateurs).
  \item Écrire un mutateur qui valide la donnée avant de l'enregistrer.
  \item Construire une classe fille par héritage, redéfinir une méthode et appeler le constructeur parent.
\end{itemize}
\textbf{Consignes} : chaque question se traite d'abord sur le papier, puis se vérifie dans
l'interpréteur Python (fichier \py{tp\_poo.py}). Le code rendu doit s'exécuter sans erreur.
\end{methode}

\begin{rappel}[Vocabulaire de référence]
Une \textbf{classe} est un modèle. Un \textbf{objet} (ou \textbf{instance}) est une réalisation
concrète de ce modèle. Un \textbf{attribut} est une donnée portée par l'objet. Une
\textbf{méthode} est une fonction définie dans la classe et appelée sur un objet. Le paramètre
\py{self} désigne l'objet sur lequel la méthode est appelée.
\end{rappel}

% ============================================================================
\section*{Partie A --- Classes, attributs et méthodes}
\addcontentsline{toc}{section}{Partie A}
\setcounter{exo}{0}

% ---------------------------------------------------------------- A1
\exo[La classe \py{Voiture}]{Lire le code suivant, puis répondre aux questions.}

\begin{lstlisting}
class Voiture:
    def __init__(self, marque, couleur):
        self.marque = marque
        self.couleur = couleur

    def afficher_details(self):
        return f"Voiture {self.marque} de couleur {self.couleur}"
\end{lstlisting}

\textbf{Questions}
\begin{enumerate}
  \item Que fait la méthode \py{__init__} ?
    \rep{C'est le \textbf{constructeur}. Il est appelé automatiquement à la création de l'objet
    et initialise les attributs d'instance \py{self.marque} et \py{self.couleur} à partir des
    valeurs passées en paramètres.}
  \item Combien d'attributs possède la classe \py{Voiture} ?
    \rep{Deux attributs d'instance : \py{marque} et \py{couleur}.}
  \item Que retourne la méthode \py{afficher_details} ?
    \rep{Une chaîne de caractères (\py{str}) construite par une f-string, par exemple
    \py{"Voiture Toyota de couleur Bleu"}. Elle \emph{retourne} la chaîne, elle ne l'affiche pas.}
\end{enumerate}

\textbf{À compléter} : créer une instance de \py{Voiture} de marque \py{"Toyota"} et de couleur
\py{"Bleu"}, appeler \py{afficher_details} et afficher le résultat.
\zonecode[3cm]
\begin{solution}
\begin{lstlisting}
ma_voiture = Voiture("Toyota", "Bleu")
print(ma_voiture.afficher_details())
# Voiture Toyota de couleur Bleu
\end{lstlisting}
\py{print} est nécessaire : sans lui, la chaîne est retournée puis perdue.
\end{solution}

% ---------------------------------------------------------------- A2
\exo[Compléter la classe \py{Chien}]{}

\begin{lstlisting}
class Chien:
    def __init__(self, nom, age):
        self.nom = nom
        self.age = age

    # Méthode aboyer qui retourne "Wouf !"
\end{lstlisting}

\textbf{Questions}
\begin{enumerate}
  \item Quels sont les attributs de la classe \py{Chien} ?
    \rep{\py{nom} et \py{age}, deux attributs d'instance.}
  \item Comment appelle-t-on la méthode \py{__init__} dans le contexte de la POO ?
    \rep{Le \textbf{constructeur} (ou initialiseur). C'est une \emph{méthode spéciale}, reconnaissable
    à son double souligné ; on ne l'appelle jamais directement, elle est déclenchée par
    \py{Chien("Rex", 3)}.}
  \item Ajouter une méthode \py{aboyer} qui retourne \py{"Wouf !"}.
  \item Créer un objet \py{Chien} nommé \py{"Rex"}, âgé de 3 ans, et le faire aboyer.
\end{enumerate}
\zonecode[4.5cm]
\begin{solution}
\begin{lstlisting}
class Chien:
    def __init__(self, nom, age):
        self.nom = nom
        self.age = age

    def aboyer(self):
        return "Wouf !"

rex = Chien("Rex", 3)
print(rex.aboyer())   # Wouf !
\end{lstlisting}
Le paramètre \py{self} est obligatoire dans la définition, même si \py{aboyer} n'utilise aucun
attribut : il reçoit l'objet sur lequel la méthode est appelée.
\end{solution}

% ---------------------------------------------------------------- A3
\exo[La classe \py{Livre}]{}

\begin{lstlisting}
class Livre:
    def __init__(self, titre, auteur):
        self.titre = titre
        self.auteur = auteur

    def description(self):
        return f"{self.titre} par {self.auteur}"
\end{lstlisting}

\textbf{Questions}
\begin{enumerate}
  \item Que fait la méthode \py{description} ?
    \rep{Elle retourne une chaîne qui associe le titre et l'auteur du livre, sous la forme
    \py{"titre par auteur"}.}
  \item Combien de paramètres la méthode \py{__init__} attend-elle, en incluant \py{self} ?
    \rep{Trois : \py{self}, \py{titre} et \py{auteur}. À l'appel, seuls deux arguments sont
    fournis : \py{self} est transmis automatiquement par Python.}
  \item Créer un objet \py{Livre} de titre \py{"Le Petit Prince"} et d'auteur
    \py{"Antoine de Saint-Exupéry"}, puis afficher sa description.
\end{enumerate}
\zonecode[3cm]
\begin{solution}
\begin{lstlisting}
livre = Livre("Le Petit Prince", "Antoine de Saint-Exupery")
print(livre.description())
# Le Petit Prince par Antoine de Saint-Exupery
\end{lstlisting}
\end{solution}

% ---------------------------------------------------------------- A4
\exo[La classe \py{Rectangle}]{}

\begin{lstlisting}
class Rectangle:
    def __init__(self, longueur, largeur):
        self.longueur = longueur
        self.largeur = largeur

    # Méthode calculer_aire qui retourne l'aire du rectangle
\end{lstlisting}

\textbf{Questions}
\begin{enumerate}
  \item Quels sont les attributs de la classe \py{Rectangle} ?
    \rep{\py{longueur} et \py{largeur}.}
  \item Quelle formule mathématique faut-il utiliser pour calculer l'aire ?
    \rep{$\text{aire} = \text{longueur} \times \text{largeur}$.}
  \item Ajouter une méthode \py{calculer_aire} qui retourne l'aire du rectangle.
  \item Créer un rectangle de longueur 5 et de largeur 3, puis afficher son aire.
\end{enumerate}
\zonecode[4.5cm]
\begin{solution}
\begin{lstlisting}
class Rectangle:
    def __init__(self, longueur, largeur):
        self.longueur = longueur
        self.largeur = largeur

    def calculer_aire(self):
        return self.longueur * self.largeur

r = Rectangle(5, 3)
print(r.calculer_aire())   # 15
\end{lstlisting}
Une méthode peut aussi être ajoutée pour le périmètre :\\
\py{return 2 * (self.longueur + self.largeur)}.
\end{solution}

% ---------------------------------------------------------------- A5
\exo[Création complète d'une classe \py{Etudiant}]{On souhaite modéliser un étudiant par son
nom et sa note.}

\begin{enumerate}
  \item Créer une classe \py{Etudiant} possédant un constructeur \py{__init__} prenant en
    paramètres \py{nom} et \py{note}, et deux attributs d'instance \py{self.nom} et \py{self.note}.
  \item Ajouter une méthode \py{a_reussi} qui retourne \py{True} si la note est supérieure ou
    égale à 10, et \py{False} sinon.
  \item Tester la classe en créant un étudiant avec une note de 8.
  \item Ajouter une méthode \py{augmenter_note} qui augmente la note de 2 points.
  \item Créer un étudiant nommé \py{"Alice"} avec une note de 15 et afficher si elle a réussi,
    grâce à la méthode \py{a_reussi}.
\end{enumerate}
\zonecode[7cm]
\begin{solution}
\begin{lstlisting}
class Etudiant:
    def __init__(self, nom, note):
        self.nom = nom
        self.note = note

    def a_reussi(self):
        return self.note >= 10

    def augmenter_note(self):
        self.note = self.note + 2

# 3. test avec une note de 8
e = Etudiant("Bob", 8)
print(e.a_reussi())          # False

# 5. Alice avec une note de 15
alice = Etudiant("Alice", 15)
print(alice.a_reussi())      # True
\end{lstlisting}
\textbf{Points d'attention.}
\py{return self.note >= 10} suffit : la comparaison produit déjà un booléen, un
\py{if ... return True else return False} est inutile.
\py{augmenter_note} \emph{modifie} l'attribut et ne retourne rien : elle agit sur l'état de l'objet.
\end{solution}

% ============================================================================
\clearpage
\section*{Partie B --- Encapsulation et héritage}
\addcontentsline{toc}{section}{Partie B}
\setcounter{exo}{0}

\begin{rappel}[Encapsulation]
Préfixer un attribut par deux soulignés (\py{self.__nom}) déclenche le \emph{name mangling} :
Python renomme l'attribut en \py{_Classe__nom}, ce qui le rend inaccessible par son nom
d'origine depuis l'extérieur. On y accède alors par un \textbf{accesseur} (\emph{getter},
\py{get_nom}) et on le modifie par un \textbf{mutateur} (\emph{setter}, \py{set_nom}), ce qui
permet de contrôler les valeurs acceptées.
\end{rappel}

\begin{rappel}[Héritage]
\py{class Chat(Animal)} signifie que \py{Chat} \textbf{hérite} de \py{Animal} : la classe fille
reçoit les attributs et méthodes de la classe mère. Elle peut en ajouter, ou
\textbf{redéfinir} une méthode existante (\emph{override}). \py{super().__init__(...)} appelle le
constructeur de la classe mère.
\end{rappel}

% ---------------------------------------------------------------- B1
\exo[Encapsulation simple avec attribut privé]{}

\begin{lstlisting}
class Personne:
    def __init__(self, nom, age):
        self.__nom = nom
        self.__age = age

    def get_nom(self):
        return self.__nom

    # Méthode get_age()
\end{lstlisting}

\textbf{Questions}
\begin{enumerate}
  \item Pourquoi les attributs \py{__nom} et \py{__age} sont-ils privés ?
    \rep{Pour empêcher leur modification directe depuis l'extérieur de la classe. L'accès passe
    par des méthodes, ce qui permet de contrôler la lecture et l'écriture et de modifier
    l'implémentation interne sans casser le code qui utilise la classe.}
  \item Comment accède-t-on à l'attribut \py{__nom} ?
    \rep{Par l'accesseur : \py{p.get_nom()}. L'écriture \py{p.__nom} provoque une
    \py{AttributeError}.}
  \item Ajouter une méthode \py{get_age} pour accéder à l'attribut \py{__age}.
  \item Créer une instance de \py{Personne} de nom \py{"Luc"} et d'âge 25, puis afficher son nom.
\end{enumerate}
\zonecode[4.5cm]
\begin{solution}
\begin{lstlisting}
class Personne:
    def __init__(self, nom, age):
        self.__nom = nom
        self.__age = age

    def get_nom(self):
        return self.__nom

    def get_age(self):
        return self.__age

p = Personne("Luc", 25)
print(p.get_nom())    # Luc
print(p.get_age())    # 25
\end{lstlisting}
\end{solution}

% ---------------------------------------------------------------- B2
\exo[Ajouter un mutateur]{}

\begin{lstlisting}
class CompteBancaire:
    def __init__(self, titulaire, solde):
        self.__titulaire = titulaire
        self.__solde = solde

    def get_solde(self):
        return self.__solde
\end{lstlisting}

\textbf{Questions}
\begin{enumerate}
  \item Que signifie l'encapsulation dans cette classe ?
    \rep{Le titulaire et le solde sont privés : aucun code extérieur ne peut écrire
    \py{compte.__solde = 1000000}. Le solde ne peut être consulté que par \py{get_solde} et
    modifié que par les méthodes que la classe met à disposition.}
  \item Que fait la méthode \py{get_solde} ?
    \rep{Elle retourne la valeur de l'attribut privé \py{__solde} : c'est un accesseur en
    lecture seule.}
  \item Créer un compte bancaire pour \py{"Emma"} avec un solde de 500 et afficher le solde.
  \item Ajouter une méthode \py{set_solde} pour modifier le solde.
\end{enumerate}
\zonecode[5cm]
\begin{solution}
\begin{lstlisting}
class CompteBancaire:
    def __init__(self, titulaire, solde):
        self.__titulaire = titulaire
        self.__solde = solde

    def get_solde(self):
        return self.__solde

    def set_solde(self, nouveau_solde):
        self.__solde = nouveau_solde

compte = CompteBancaire("Emma", 500)
print(compte.get_solde())   # 500
compte.set_solde(750)
print(compte.get_solde())   # 750
\end{lstlisting}
\textbf{Prolongement.} Un mutateur brut n'apporte pas de sécurité. Une version plus réaliste
refuse un solde négatif, ou expose plutôt des opérations métier :
\py{deposer(montant)} et \py{retirer(montant)} avec vérification de la provision.
\end{solution}

% ---------------------------------------------------------------- B3
\exo[Encapsulation avec validation]{}

\begin{lstlisting}
class Etudiant:
    def __init__(self, nom, note):
        self.__nom = nom
        self.__note = note

    def get_note(self):
        return self.__note

    # Méthode set_note qui accepte uniquement des notes entre 0 et 20
\end{lstlisting}

\textbf{Questions}
\begin{enumerate}
  \item Pourquoi utiliser un mutateur pour la note ?
    \rep{Pour garantir un \emph{invariant} : une note reste comprise entre 0 et 20. Le mutateur
    est le seul point d'entrée en écriture, il peut donc vérifier la valeur avant de
    l'enregistrer et refuser les valeurs incohérentes.}
  \item Que se passe-t-il si on essaie d'accéder directement à \py{__note} ?
    \rep{Python lève une \py{AttributeError} : à cause du \emph{name mangling}, l'attribut
    s'appelle en réalité \py{_Etudiant__note} et le nom \py{__note} n'existe pas depuis
    l'extérieur.}
  \item Ajouter une méthode \py{set_note} qui n'accepte que des notes entre 0 et 20.
  \item Créer un étudiant nommé \py{"Paul"} avec une note de 15, puis essayer de changer la note
    à 25.
\end{enumerate}
\zonecode[5.5cm]
\begin{solution}
\begin{lstlisting}
class Etudiant:
    def __init__(self, nom, note):
        self.__nom = nom
        self.__note = note

    def get_note(self):
        return self.__note

    def set_note(self, nouvelle_note):
        if 0 <= nouvelle_note <= 20:
            self.__note = nouvelle_note
        else:
            print("Note invalide : elle doit etre comprise entre 0 et 20.")

paul = Etudiant("Paul", 15)
print(paul.get_note())   # 15
paul.set_note(25)        # Note invalide : ...
print(paul.get_note())   # 15 : la note n'a pas change
\end{lstlisting}
\textbf{Variante professionnelle.} Au lieu d'un \py{print}, lever une exception rend l'erreur
impossible à ignorer : \py{raise ValueError("note hors intervalle")}.
\end{solution}

% ---------------------------------------------------------------- B4
\exo[Héritage simple]{}

\begin{lstlisting}
class Animal:
    def __init__(self, nom):
        self.__nom = nom

    def get_nom(self):
        return self.__nom

    def faire_son(self):
        return "Son generique"

class Chat(Animal):
    # Méthode faire_son pour retourner "Miaou"
\end{lstlisting}

\textbf{Questions}
\begin{enumerate}
  \item Quelle classe est la classe parent ?
    \rep{\py{Animal}. \py{Chat} est la classe fille (ou classe dérivée).}
  \item Que fait l'héritage dans ce cas ?
    \rep{\py{Chat} récupère sans les réécrire le constructeur \py{__init__}, l'attribut privé
    \py{__nom} et la méthode \py{get_nom} d'\py{Animal}. Elle peut en plus redéfinir
    \py{faire_son} pour lui donner son comportement propre.}
  \item Redéfinir la méthode \py{faire_son} pour qu'elle retourne \py{"Miaou"}.
  \item Créer un \py{Chat} nommé \py{"Minou"} et afficher son son.
\end{enumerate}
\zonecode[4.5cm]
\begin{solution}
\begin{lstlisting}
class Chat(Animal):
    def faire_son(self):
        return "Miaou"

minou = Chat("Minou")
print(minou.get_nom())     # Minou  (methode heritee)
print(minou.faire_son())   # Miaou  (methode redefinie)
\end{lstlisting}
\py{Chat} ne définit pas de \py{__init__} : celui d'\py{Animal} est utilisé tel quel.
Appeler \py{faire_son()} sur un \py{Chat} exécute la version de \py{Chat} : c'est le
\textbf{polymorphisme}.
\end{solution}

% ---------------------------------------------------------------- B5
\exo[Héritage avec nouvel attribut]{}

\begin{lstlisting}
class Vehicule:
    def __init__(self, marque):
        self.__marque = marque

    def get_marque(self):
        return self.__marque

class Voiture(Vehicule):
    def __init__(self, marque, couleur):
        super().__init__(marque)
        self.__couleur = couleur

    # À compléter : ajoutez une méthode get_couleur
\end{lstlisting}

\textbf{Questions}
\begin{enumerate}
  \item Que fait \py{super().__init__(marque)} ?
    \rep{Il appelle le constructeur de la classe mère \py{Vehicule} pour initialiser la partie
    héritée de l'objet, ici l'attribut privé \py{__marque}. Sans cet appel, \py{get_marque}
    échouerait car l'attribut n'aurait jamais été créé.}
  \item Combien d'attributs privés possède la classe \py{Voiture} ?
    \rep{Deux au total dans l'objet : \py{_Vehicule__marque}, hérité et initialisé par
    \py{super()}, et \py{_Voiture__couleur}, propre à \py{Voiture}. La classe \py{Voiture} n'en
    déclare qu'un seul elle-même.}
  \item À compléter : ajouter une méthode \py{get_couleur}.
\end{enumerate}
\zonecode[4cm]
\begin{solution}
\begin{lstlisting}
class Voiture(Vehicule):
    def __init__(self, marque, couleur):
        super().__init__(marque)
        self.__couleur = couleur

    def get_couleur(self):
        return self.__couleur

v = Voiture("Renault", "Bleu")
print(v.get_marque())    # Renault
print(v.get_couleur())   # Bleu
\end{lstlisting}
\begin{piege}
\py{__marque} est privé \emph{à \py{Vehicule}} : à cause du \emph{name mangling},
une méthode de \py{Voiture} ne peut pas écrire \py{self.__marque}. Elle doit passer par
\py{self.get_marque()}. C'est exactement le rôle de l'accesseur hérité.
\end{piege}
\end{solution}

% ---------------------------------------------------------------- B6
\exo[Encapsulation et mutateur avec validation]{}

\begin{lstlisting}
class Livre:
    def __init__(self, titre, prix):
        self.__titre = titre
        self.__prix = prix

    def get_prix(self):
        return self.__prix

    # À compléter : ajoutez une méthode set_prix qui refuse les prix négatifs
\end{lstlisting}

\textbf{Questions}
\begin{enumerate}
  \item Pourquoi le prix est-il un attribut privé ?
    \rep{Parce qu'un prix doit respecter une règle métier : il ne peut pas être négatif. Le
    rendre privé interdit l'écriture directe et oblige à passer par une méthode qui applique
    cette règle.}
  \item Que fait un mutateur avec validation ?
    \rep{Il vérifie la valeur reçue avant de l'affecter à l'attribut. Si la valeur est valide,
    il met l'attribut à jour ; sinon il refuse la modification (message d'erreur ou exception)
    et l'objet conserve son état précédent, qui reste cohérent.}
  \item Créer un \py{Livre} intitulé \py{"Python 101"} avec un prix de 30, puis essayer de
    définir un prix de $-10$.
\end{enumerate}
\zonecode[5.5cm]
\begin{solution}
\begin{lstlisting}
class Livre:
    def __init__(self, titre, prix):
        self.__titre = titre
        self.__prix = prix

    def get_prix(self):
        return self.__prix

    def set_prix(self, nouveau_prix):
        if nouveau_prix >= 0:
            self.__prix = nouveau_prix
        else:
            print("Prix invalide : un prix ne peut pas etre negatif.")

livre = Livre("Python 101", 30)
print(livre.get_prix())   # 30
livre.set_prix(-10)       # Prix invalide : ...
print(livre.get_prix())   # 30 : le prix est inchange
\end{lstlisting}
\textbf{Remarque.} Le constructeur devrait appliquer la même validation, par exemple en
appelant \py{self.set_prix(prix)} au lieu d'affecter directement : sinon
\py{Livre("X", -10)} crée un objet invalide dès le départ.
\end{solution}

% ============================================================================
\clearpage
\section*{Synthèse}

\begin{center}
\begin{tabularx}{\textwidth}{@{}>{\bfseries}p{4.3cm} X@{}}
\toprule
Notion & Ce qu'il faut retenir \\
\midrule
Classe & Modèle décrivant des données et des comportements : \py{class Voiture:} \\
Objet / instance & Réalisation concrète de la classe : \py{v = Voiture("Toyota", "Bleu")} \\
Constructeur & \py{__init__}, appelé automatiquement, initialise les attributs \\
\py{self} & Référence à l'objet courant ; premier paramètre de toute méthode d'instance \\
Attribut d'instance & Donnée propre à l'objet : \py{self.marque} \\
Méthode & Fonction définie dans la classe, appelée par \py{objet.methode()} \\
Attribut privé & \py{self.__prix} : inaccessible depuis l'extérieur (\emph{name mangling}) \\
Accesseur (\emph{getter}) & \py{get_prix()} : lit un attribut privé \\
Mutateur (\emph{setter}) & \py{set_prix(v)} : écrit un attribut privé après validation \\
Héritage & \py{class Voiture(Vehicule)} : la fille reçoit les membres de la mère \\
\py{super()} & Accède à la classe mère, notamment \py{super().__init__(...)} \\
Redéfinition & La fille réécrit une méthode de la mère (\py{faire_son}) \\
\bottomrule
\end{tabularx}
\end{center}

\begin{piege}
\textbf{Erreurs les plus fréquentes.}
oublier \py{self} dans la définition d'une méthode ;
confondre \py{return} (retourne une valeur) et \py{print} (affiche) ;
écrire \py{objet.__attribut} depuis l'extérieur ;
oublier \py{super().__init__(...)} dans le constructeur d'une classe fille ;
appeler une méthode sans les parenthèses (\py{v.afficher_details} retourne la méthode,
pas son résultat).
\end{piege}

\end{document}
"
\documentclass[11pt,a4paper]{article}

\usepackage[T1]{fontenc}
\usepackage[utf8]{inputenc}
\usepackage{lmodern}
\usepackage[french]{babel}
\usepackage[a4paper,margin=1.8cm,bottom=2.2cm]{geometry}
\usepackage[table]{xcolor}
\usepackage{booktabs,tabularx,array}
\usepackage{enumitem}
\usepackage{needspace}
\usepackage{fancyhdr}
\usepackage{amssymb,amsmath}
\usepackage[most]{tcolorbox}
\usepackage{listings}
\usepackage[strings]{underscore}  % permet d'écrire __init__ sans échappement
\usepackage{hyperref}
\hypersetup{hidelinks}

\providecommand{\modecorrige}{0}
\newif\ifcorrige
\ifnum\modecorrige=1 \corrigetrue \else \corrigefalse \fi

% Palette alignée sur les autres supports du dépôt
\definecolor{ipred}{HTML}{C0392B}
\definecolor{cdcbar}{HTML}{A67C3D}
\definecolor{cdcbg}{HTML}{F4F1EA}
\definecolor{qbar}{HTML}{4A6B5D}
\definecolor{qbg}{HTML}{EDF1EF}
\definecolor{corr}{HTML}{1F5FA8}
\definecolor{corrbg}{HTML}{E8F0FA}
\definecolor{grisfoot}{HTML}{828282}
\definecolor{codebg}{HTML}{F7F7F4}
\definecolor{codekw}{HTML}{8E44AD}
\definecolor{codestr}{HTML}{1F7A4D}
\definecolor{codecom}{HTML}{8A8A8A}

% ---------------------------------------------------------------- code Python
\lstdefinestyle{py}{
  language=Python,
  basicstyle=\ttfamily\small,
  keywordstyle=\color{codekw}\bfseries,
  stringstyle=\color{codestr},
  commentstyle=\color{codecom}\itshape,
  showstringspaces=false,
  backgroundcolor=\color{codebg},
  frame=single,
  rulecolor=\color{black!25},
  framesep=5pt,
  xleftmargin=4pt,
  xrightmargin=4pt,
  columns=fullflexible,
  keepspaces=true,
  breaklines=true,
  extendedchars=true,
  literate={é}{{\'e}}1 {è}{{\`e}}1 {ê}{{\^e}}1 {à}{{\`a}}1 {ç}{{\c c}}1
           {ù}{{\`u}}1 {û}{{\^u}}1 {î}{{\^\i}}1 {ô}{{\^o}}1 {É}{{\'E}}1 {À}{{\`A}}1 {È}{{\`E}}1 {Ç}{{\c C}}1
           {«}{{\guillemotleft}}1 {»}{{\guillemotright}}1,
  morekeywords={self,True,False,None},
}
\lstset{style=py}
\newcommand{\py}[1]{\texttt{#1}}

% --------------------------------------------------------------- encadrés
\newtcolorbox{methode}[1][]{enhanced,breakable,colback=cdcbg,colframe=cdcbar,boxrule=0pt,
  leftrule=3pt,arc=0pt,outer arc=0pt,left=6pt,right=6pt,top=4pt,bottom=4pt,
  title={#1},coltitle=cdcbar,fonttitle=\bfseries,detach title,
  before upper={\tcbtitle\par\smallskip}}
\newtcolorbox{rappel}[1][]{enhanced,breakable,colback=qbg,colframe=qbar,boxrule=0pt,
  leftrule=3pt,arc=0pt,outer arc=0pt,left=6pt,right=6pt,top=4pt,bottom=4pt,
  title={#1},coltitle=qbar,fonttitle=\bfseries,attach title to upper={\ :\ }}
\newtcolorbox{piege}{enhanced,breakable,colback=red!4,colframe=ipred,boxrule=0pt,
  leftrule=3pt,arc=0pt,outer arc=0pt,left=6pt,right=6pt,top=4pt,bottom=4pt}

% --------------------------------------------------------------- exercices
\newcounter{exo}[section]
\renewcommand{\theexo}{\arabic{exo}}
\newtcolorbox{quest}{enhanced,breakable,colback=qbg,colframe=qbar,boxrule=0pt,
  leftrule=3pt,arc=0pt,outer arc=0pt,left=6pt,right=6pt,top=4pt,bottom=4pt}
\newcommand{\exo}[2][]{% [sous-titre] {énoncé}
  \Needspace{5cm}\stepcounter{exo}%
  \begin{quest}\textbf{Exercice \theexo\ifx\relax#1\relax\else\ — #1\fi.}
  \ifx\relax#2\relax\else\ #2\fi\end{quest}}

% Réponse courte : pointillés dans le sujet, texte bleu dans le corrigé
\newcommand{\rep}[1]{%
  \ifcorrige
    \par\smallskip{\color{corr}\small#1}\par\smallskip
  \else
    \par\nobreak\vspace{1mm}\noindent\color{black!55}\dotfill\par
    \noindent\color{black!55}\dotfill\par\color{black}\vspace{1mm}
  \fi}

% Zone de code vierge (sujet uniquement) de hauteur donnée
\newcommand{\zonecode}[1][4cm]{%
  \ifcorrige\else
    \par\smallskip\noindent
    \begin{tcolorbox}[colback=white,colframe=black!45,boxrule=0.5pt,arc=1pt,height=#1]\end{tcolorbox}%
  \fi}

% Bloc visible seulement dans le corrigé
\newenvironment{solution}
  {\ifcorrige
     \begin{tcolorbox}[enhanced,breakable,colback=corrbg,colframe=corr,boxrule=0.6pt,arc=1pt,
       left=5pt,right=5pt,top=4pt,bottom=4pt,
       title={Corrigé},coltitle=corr,fonttitle=\bfseries\small,detach title,
       before upper={\tcbtitle\par\smallskip},fontupper=\small]%
   \else
     \setbox0\vbox\bgroup
   \fi}
  {\ifcorrige\end{tcolorbox}\else\egroup\fi}

\pagestyle{fancy}
\fancyhf{}
\renewcommand{\headrulewidth}{0pt}
\fancyfoot[L]{\small\itshape\color{grisfoot}TP POO Python\ifcorrige\ — CORRIGÉ\fi}
\fancyfoot[C]{\small\color{grisfoot}BTS CIEL}
\fancyfoot[R]{\small\color{grisfoot}\thepage}

\setlist{itemsep=2pt,topsep=3pt}
\setlength{\parindent}{0pt}
\setlength{\parskip}{4pt}
\renewcommand{\arraystretch}{1.3}

\begin{document}

\begin{center}
{\LARGE\bfseries Programmation orientée objet en Python\par}
\smallskip
{\large Classes, attributs, méthodes --- Encapsulation et héritage\par}
\medskip
BTS CIEL --- Travaux pratiques
\ifcorrige\par\textcolor{corr}{\textbf{Version enseignant --- corrigés intégrés}}\fi
\end{center}

\medskip
\noindent\begin{tabularx}{\textwidth}{@{}>{\bfseries}l X@{}}
Nom, prénom : & \ifcorrige\textcolor{corr}{corrigé}\else\dotfill\fi \\
Classe : & \ifcorrige\textcolor{corr}{---}\else\dotfill\fi \\
Date : & \ifcorrige\textcolor{corr}{---}\else\dotfill\fi \\
\end{tabularx}

\begin{methode}[Objectifs du TP]
\begin{itemize}
  \item Lire une classe Python et identifier son constructeur, ses attributs et ses méthodes.
  \item Instancier un objet et appeler ses méthodes.
  \item Compléter une classe existante en ajoutant une méthode.
  \item Protéger les données d'un objet par l'encapsulation (attributs privés, accesseurs, mutateurs).
  \item Écrire un mutateur qui valide la donnée avant de l'enregistrer.
  \item Construire une classe fille par héritage, redéfinir une méthode et appeler le constructeur parent.
\end{itemize}
\textbf{Consignes} : chaque question se traite d'abord sur le papier, puis se vérifie dans
l'interpréteur Python (fichier \py{tp\_poo.py}). Le code rendu doit s'exécuter sans erreur.
\end{methode}

\begin{rappel}[Vocabulaire de référence]
Une \textbf{classe} est un modèle. Un \textbf{objet} (ou \textbf{instance}) est une réalisation
concrète de ce modèle. Un \textbf{attribut} est une donnée portée par l'objet. Une
\textbf{méthode} est une fonction définie dans la classe et appelée sur un objet. Le paramètre
\py{self} désigne l'objet sur lequel la méthode est appelée.
\end{rappel}

% ============================================================================
\section*{Partie A --- Classes, attributs et méthodes}
\addcontentsline{toc}{section}{Partie A}
\setcounter{exo}{0}

% ---------------------------------------------------------------- A1
\exo[La classe \py{Voiture}]{Lire le code suivant, puis répondre aux questions.}

\begin{lstlisting}
class Voiture:
    def __init__(self, marque, couleur):
        self.marque = marque
        self.couleur = couleur

    def afficher_details(self):
        return f"Voiture {self.marque} de couleur {self.couleur}"
\end{lstlisting}

\textbf{Questions}
\begin{enumerate}
  \item Que fait la méthode \py{__init__} ?
    \rep{C'est le \textbf{constructeur}. Il est appelé automatiquement à la création de l'objet
    et initialise les attributs d'instance \py{self.marque} et \py{self.couleur} à partir des
    valeurs passées en paramètres.}
  \item Combien d'attributs possède la classe \py{Voiture} ?
    \rep{Deux attributs d'instance : \py{marque} et \py{couleur}.}
  \item Que retourne la méthode \py{afficher_details} ?
    \rep{Une chaîne de caractères (\py{str}) construite par une f-string, par exemple
    \py{"Voiture Toyota de couleur Bleu"}. Elle \emph{retourne} la chaîne, elle ne l'affiche pas.}
\end{enumerate}

\textbf{À compléter} : créer une instance de \py{Voiture} de marque \py{"Toyota"} et de couleur
\py{"Bleu"}, appeler \py{afficher_details} et afficher le résultat.
\zonecode[3cm]
\begin{solution}
\begin{lstlisting}
ma_voiture = Voiture("Toyota", "Bleu")
print(ma_voiture.afficher_details())
# Voiture Toyota de couleur Bleu
\end{lstlisting}
\py{print} est nécessaire : sans lui, la chaîne est retournée puis perdue.
\end{solution}

% ---------------------------------------------------------------- A2
\exo[Compléter la classe \py{Chien}]{}

\begin{lstlisting}
class Chien:
    def __init__(self, nom, age):
        self.nom = nom
        self.age = age

    # Méthode aboyer qui retourne "Wouf !"
\end{lstlisting}

\textbf{Questions}
\begin{enumerate}
  \item Quels sont les attributs de la classe \py{Chien} ?
    \rep{\py{nom} et \py{age}, deux attributs d'instance.}
  \item Comment appelle-t-on la méthode \py{__init__} dans le contexte de la POO ?
    \rep{Le \textbf{constructeur} (ou initialiseur). C'est une \emph{méthode spéciale}, reconnaissable
    à son double souligné ; on ne l'appelle jamais directement, elle est déclenchée par
    \py{Chien("Rex", 3)}.}
  \item Ajouter une méthode \py{aboyer} qui retourne \py{"Wouf !"}.
  \item Créer un objet \py{Chien} nommé \py{"Rex"}, âgé de 3 ans, et le faire aboyer.
\end{enumerate}
\zonecode[4.5cm]
\begin{solution}
\begin{lstlisting}
class Chien:
    def __init__(self, nom, age):
        self.nom = nom
        self.age = age

    def aboyer(self):
        return "Wouf !"

rex = Chien("Rex", 3)
print(rex.aboyer())   # Wouf !
\end{lstlisting}
Le paramètre \py{self} est obligatoire dans la définition, même si \py{aboyer} n'utilise aucun
attribut : il reçoit l'objet sur lequel la méthode est appelée.
\end{solution}

% ---------------------------------------------------------------- A3
\exo[La classe \py{Livre}]{}

\begin{lstlisting}
class Livre:
    def __init__(self, titre, auteur):
        self.titre = titre
        self.auteur = auteur

    def description(self):
        return f"{self.titre} par {self.auteur}"
\end{lstlisting}

\textbf{Questions}
\begin{enumerate}
  \item Que fait la méthode \py{description} ?
    \rep{Elle retourne une chaîne qui associe le titre et l'auteur du livre, sous la forme
    \py{"titre par auteur"}.}
  \item Combien de paramètres la méthode \py{__init__} attend-elle, en incluant \py{self} ?
    \rep{Trois : \py{self}, \py{titre} et \py{auteur}. À l'appel, seuls deux arguments sont
    fournis : \py{self} est transmis automatiquement par Python.}
  \item Créer un objet \py{Livre} de titre \py{"Le Petit Prince"} et d'auteur
    \py{"Antoine de Saint-Exupéry"}, puis afficher sa description.
\end{enumerate}
\zonecode[3cm]
\begin{solution}
\begin{lstlisting}
livre = Livre("Le Petit Prince", "Antoine de Saint-Exupery")
print(livre.description())
# Le Petit Prince par Antoine de Saint-Exupery
\end{lstlisting}
\end{solution}

% ---------------------------------------------------------------- A4
\exo[La classe \py{Rectangle}]{}

\begin{lstlisting}
class Rectangle:
    def __init__(self, longueur, largeur):
        self.longueur = longueur
        self.largeur = largeur

    # Méthode calculer_aire qui retourne l'aire du rectangle
\end{lstlisting}

\textbf{Questions}
\begin{enumerate}
  \item Quels sont les attributs de la classe \py{Rectangle} ?
    \rep{\py{longueur} et \py{largeur}.}
  \item Quelle formule mathématique faut-il utiliser pour calculer l'aire ?
    \rep{$\text{aire} = \text{longueur} \times \text{largeur}$.}
  \item Ajouter une méthode \py{calculer_aire} qui retourne l'aire du rectangle.
  \item Créer un rectangle de longueur 5 et de largeur 3, puis afficher son aire.
\end{enumerate}
\zonecode[4.5cm]
\begin{solution}
\begin{lstlisting}
class Rectangle:
    def __init__(self, longueur, largeur):
        self.longueur = longueur
        self.largeur = largeur

    def calculer_aire(self):
        return self.longueur * self.largeur

r = Rectangle(5, 3)
print(r.calculer_aire())   # 15
\end{lstlisting}
Une méthode peut aussi être ajoutée pour le périmètre :\\
\py{return 2 * (self.longueur + self.largeur)}.
\end{solution}

% ---------------------------------------------------------------- A5
\exo[Création complète d'une classe \py{Etudiant}]{On souhaite modéliser un étudiant par son
nom et sa note.}

\begin{enumerate}
  \item Créer une classe \py{Etudiant} possédant un constructeur \py{__init__} prenant en
    paramètres \py{nom} et \py{note}, et deux attributs d'instance \py{self.nom} et \py{self.note}.
  \item Ajouter une méthode \py{a_reussi} qui retourne \py{True} si la note est supérieure ou
    égale à 10, et \py{False} sinon.
  \item Tester la classe en créant un étudiant avec une note de 8.
  \item Ajouter une méthode \py{augmenter_note} qui augmente la note de 2 points.
  \item Créer un étudiant nommé \py{"Alice"} avec une note de 15 et afficher si elle a réussi,
    grâce à la méthode \py{a_reussi}.
\end{enumerate}
\zonecode[7cm]
\begin{solution}
\begin{lstlisting}
class Etudiant:
    def __init__(self, nom, note):
        self.nom = nom
        self.note = note

    def a_reussi(self):
        return self.note >= 10

    def augmenter_note(self):
        self.note = self.note + 2

# 3. test avec une note de 8
e = Etudiant("Bob", 8)
print(e.a_reussi())          # False

# 5. Alice avec une note de 15
alice = Etudiant("Alice", 15)
print(alice.a_reussi())      # True
\end{lstlisting}
\textbf{Points d'attention.}
\py{return self.note >= 10} suffit : la comparaison produit déjà un booléen, un
\py{if ... return True else return False} est inutile.
\py{augmenter_note} \emph{modifie} l'attribut et ne retourne rien : elle agit sur l'état de l'objet.
\end{solution}

% ============================================================================
\clearpage
\section*{Partie B --- Encapsulation et héritage}
\addcontentsline{toc}{section}{Partie B}
\setcounter{exo}{0}

\begin{rappel}[Encapsulation]
Préfixer un attribut par deux soulignés (\py{self.__nom}) déclenche le \emph{name mangling} :
Python renomme l'attribut en \py{_Classe__nom}, ce qui le rend inaccessible par son nom
d'origine depuis l'extérieur. On y accède alors par un \textbf{accesseur} (\emph{getter},
\py{get_nom}) et on le modifie par un \textbf{mutateur} (\emph{setter}, \py{set_nom}), ce qui
permet de contrôler les valeurs acceptées.
\end{rappel}

\begin{rappel}[Héritage]
\py{class Chat(Animal)} signifie que \py{Chat} \textbf{hérite} de \py{Animal} : la classe fille
reçoit les attributs et méthodes de la classe mère. Elle peut en ajouter, ou
\textbf{redéfinir} une méthode existante (\emph{override}). \py{super().__init__(...)} appelle le
constructeur de la classe mère.
\end{rappel}

% ---------------------------------------------------------------- B1
\exo[Encapsulation simple avec attribut privé]{}

\begin{lstlisting}
class Personne:
    def __init__(self, nom, age):
        self.__nom = nom
        self.__age = age

    def get_nom(self):
        return self.__nom

    # Méthode get_age()
\end{lstlisting}

\textbf{Questions}
\begin{enumerate}
  \item Pourquoi les attributs \py{__nom} et \py{__age} sont-ils privés ?
    \rep{Pour empêcher leur modification directe depuis l'extérieur de la classe. L'accès passe
    par des méthodes, ce qui permet de contrôler la lecture et l'écriture et de modifier
    l'implémentation interne sans casser le code qui utilise la classe.}
  \item Comment accède-t-on à l'attribut \py{__nom} ?
    \rep{Par l'accesseur : \py{p.get_nom()}. L'écriture \py{p.__nom} provoque une
    \py{AttributeError}.}
  \item Ajouter une méthode \py{get_age} pour accéder à l'attribut \py{__age}.
  \item Créer une instance de \py{Personne} de nom \py{"Luc"} et d'âge 25, puis afficher son nom.
\end{enumerate}
\zonecode[4.5cm]
\begin{solution}
\begin{lstlisting}
class Personne:
    def __init__(self, nom, age):
        self.__nom = nom
        self.__age = age

    def get_nom(self):
        return self.__nom

    def get_age(self):
        return self.__age

p = Personne("Luc", 25)
print(p.get_nom())    # Luc
print(p.get_age())    # 25
\end{lstlisting}
\end{solution}

% ---------------------------------------------------------------- B2
\exo[Ajouter un mutateur]{}

\begin{lstlisting}
class CompteBancaire:
    def __init__(self, titulaire, solde):
        self.__titulaire = titulaire
        self.__solde = solde

    def get_solde(self):
        return self.__solde
\end{lstlisting}

\textbf{Questions}
\begin{enumerate}
  \item Que signifie l'encapsulation dans cette classe ?
    \rep{Le titulaire et le solde sont privés : aucun code extérieur ne peut écrire
    \py{compte.__solde = 1000000}. Le solde ne peut être consulté que par \py{get_solde} et
    modifié que par les méthodes que la classe met à disposition.}
  \item Que fait la méthode \py{get_solde} ?
    \rep{Elle retourne la valeur de l'attribut privé \py{__solde} : c'est un accesseur en
    lecture seule.}
  \item Créer un compte bancaire pour \py{"Emma"} avec un solde de 500 et afficher le solde.
  \item Ajouter une méthode \py{set_solde} pour modifier le solde.
\end{enumerate}
\zonecode[5cm]
\begin{solution}
\begin{lstlisting}
class CompteBancaire:
    def __init__(self, titulaire, solde):
        self.__titulaire = titulaire
        self.__solde = solde

    def get_solde(self):
        return self.__solde

    def set_solde(self, nouveau_solde):
        self.__solde = nouveau_solde

compte = CompteBancaire("Emma", 500)
print(compte.get_solde())   # 500
compte.set_solde(750)
print(compte.get_solde())   # 750
\end{lstlisting}
\textbf{Prolongement.} Un mutateur brut n'apporte pas de sécurité. Une version plus réaliste
refuse un solde négatif, ou expose plutôt des opérations métier :
\py{deposer(montant)} et \py{retirer(montant)} avec vérification de la provision.
\end{solution}

% ---------------------------------------------------------------- B3
\exo[Encapsulation avec validation]{}

\begin{lstlisting}
class Etudiant:
    def __init__(self, nom, note):
        self.__nom = nom
        self.__note = note

    def get_note(self):
        return self.__note

    # Méthode set_note qui accepte uniquement des notes entre 0 et 20
\end{lstlisting}

\textbf{Questions}
\begin{enumerate}
  \item Pourquoi utiliser un mutateur pour la note ?
    \rep{Pour garantir un \emph{invariant} : une note reste comprise entre 0 et 20. Le mutateur
    est le seul point d'entrée en écriture, il peut donc vérifier la valeur avant de
    l'enregistrer et refuser les valeurs incohérentes.}
  \item Que se passe-t-il si on essaie d'accéder directement à \py{__note} ?
    \rep{Python lève une \py{AttributeError} : à cause du \emph{name mangling}, l'attribut
    s'appelle en réalité \py{_Etudiant__note} et le nom \py{__note} n'existe pas depuis
    l'extérieur.}
  \item Ajouter une méthode \py{set_note} qui n'accepte que des notes entre 0 et 20.
  \item Créer un étudiant nommé \py{"Paul"} avec une note de 15, puis essayer de changer la note
    à 25.
\end{enumerate}
\zonecode[5.5cm]
\begin{solution}
\begin{lstlisting}
class Etudiant:
    def __init__(self, nom, note):
        self.__nom = nom
        self.__note = note

    def get_note(self):
        return self.__note

    def set_note(self, nouvelle_note):
        if 0 <= nouvelle_note <= 20:
            self.__note = nouvelle_note
        else:
            print("Note invalide : elle doit etre comprise entre 0 et 20.")

paul = Etudiant("Paul", 15)
print(paul.get_note())   # 15
paul.set_note(25)        # Note invalide : ...
print(paul.get_note())   # 15 : la note n'a pas change
\end{lstlisting}
\textbf{Variante professionnelle.} Au lieu d'un \py{print}, lever une exception rend l'erreur
impossible à ignorer : \py{raise ValueError("note hors intervalle")}.
\end{solution}

% ---------------------------------------------------------------- B4
\exo[Héritage simple]{}

\begin{lstlisting}
class Animal:
    def __init__(self, nom):
        self.__nom = nom

    def get_nom(self):
        return self.__nom

    def faire_son(self):
        return "Son generique"

class Chat(Animal):
    # Méthode faire_son pour retourner "Miaou"
\end{lstlisting}

\textbf{Questions}
\begin{enumerate}
  \item Quelle classe est la classe parent ?
    \rep{\py{Animal}. \py{Chat} est la classe fille (ou classe dérivée).}
  \item Que fait l'héritage dans ce cas ?
    \rep{\py{Chat} récupère sans les réécrire le constructeur \py{__init__}, l'attribut privé
    \py{__nom} et la méthode \py{get_nom} d'\py{Animal}. Elle peut en plus redéfinir
    \py{faire_son} pour lui donner son comportement propre.}
  \item Redéfinir la méthode \py{faire_son} pour qu'elle retourne \py{"Miaou"}.
  \item Créer un \py{Chat} nommé \py{"Minou"} et afficher son son.
\end{enumerate}
\zonecode[4.5cm]
\begin{solution}
\begin{lstlisting}
class Chat(Animal):
    def faire_son(self):
        return "Miaou"

minou = Chat("Minou")
print(minou.get_nom())     # Minou  (methode heritee)
print(minou.faire_son())   # Miaou  (methode redefinie)
\end{lstlisting}
\py{Chat} ne définit pas de \py{__init__} : celui d'\py{Animal} est utilisé tel quel.
Appeler \py{faire_son()} sur un \py{Chat} exécute la version de \py{Chat} : c'est le
\textbf{polymorphisme}.
\end{solution}

% ---------------------------------------------------------------- B5
\exo[Héritage avec nouvel attribut]{}

\begin{lstlisting}
class Vehicule:
    def __init__(self, marque):
        self.__marque = marque

    def get_marque(self):
        return self.__marque

class Voiture(Vehicule):
    def __init__(self, marque, couleur):
        super().__init__(marque)
        self.__couleur = couleur

    # À compléter : ajoutez une méthode get_couleur
\end{lstlisting}

\textbf{Questions}
\begin{enumerate}
  \item Que fait \py{super().__init__(marque)} ?
    \rep{Il appelle le constructeur de la classe mère \py{Vehicule} pour initialiser la partie
    héritée de l'objet, ici l'attribut privé \py{__marque}. Sans cet appel, \py{get_marque}
    échouerait car l'attribut n'aurait jamais été créé.}
  \item Combien d'attributs privés possède la classe \py{Voiture} ?
    \rep{Deux au total dans l'objet : \py{_Vehicule__marque}, hérité et initialisé par
    \py{super()}, et \py{_Voiture__couleur}, propre à \py{Voiture}. La classe \py{Voiture} n'en
    déclare qu'un seul elle-même.}
  \item À compléter : ajouter une méthode \py{get_couleur}.
\end{enumerate}
\zonecode[4cm]
\begin{solution}
\begin{lstlisting}
class Voiture(Vehicule):
    def __init__(self, marque, couleur):
        super().__init__(marque)
        self.__couleur = couleur

    def get_couleur(self):
        return self.__couleur

v = Voiture("Renault", "Bleu")
print(v.get_marque())    # Renault
print(v.get_couleur())   # Bleu
\end{lstlisting}
\begin{piege}
\py{__marque} est privé \emph{à \py{Vehicule}} : à cause du \emph{name mangling},
une méthode de \py{Voiture} ne peut pas écrire \py{self.__marque}. Elle doit passer par
\py{self.get_marque()}. C'est exactement le rôle de l'accesseur hérité.
\end{piege}
\end{solution}

% ---------------------------------------------------------------- B6
\exo[Encapsulation et mutateur avec validation]{}

\begin{lstlisting}
class Livre:
    def __init__(self, titre, prix):
        self.__titre = titre
        self.__prix = prix

    def get_prix(self):
        return self.__prix

    # À compléter : ajoutez une méthode set_prix qui refuse les prix négatifs
\end{lstlisting}

\textbf{Questions}
\begin{enumerate}
  \item Pourquoi le prix est-il un attribut privé ?
    \rep{Parce qu'un prix doit respecter une règle métier : il ne peut pas être négatif. Le
    rendre privé interdit l'écriture directe et oblige à passer par une méthode qui applique
    cette règle.}
  \item Que fait un mutateur avec validation ?
    \rep{Il vérifie la valeur reçue avant de l'affecter à l'attribut. Si la valeur est valide,
    il met l'attribut à jour ; sinon il refuse la modification (message d'erreur ou exception)
    et l'objet conserve son état précédent, qui reste cohérent.}
  \item Créer un \py{Livre} intitulé \py{"Python 101"} avec un prix de 30, puis essayer de
    définir un prix de $-10$.
\end{enumerate}
\zonecode[5.5cm]
\begin{solution}
\begin{lstlisting}
class Livre:
    def __init__(self, titre, prix):
        self.__titre = titre
        self.__prix = prix

    def get_prix(self):
        return self.__prix

    def set_prix(self, nouveau_prix):
        if nouveau_prix >= 0:
            self.__prix = nouveau_prix
        else:
            print("Prix invalide : un prix ne peut pas etre negatif.")

livre = Livre("Python 101", 30)
print(livre.get_prix())   # 30
livre.set_prix(-10)       # Prix invalide : ...
print(livre.get_prix())   # 30 : le prix est inchange
\end{lstlisting}
\textbf{Remarque.} Le constructeur devrait appliquer la même validation, par exemple en
appelant \py{self.set_prix(prix)} au lieu d'affecter directement : sinon
\py{Livre("X", -10)} crée un objet invalide dès le départ.
\end{solution}

% ============================================================================
\clearpage
\section*{Synthèse}

\begin{center}
\begin{tabularx}{\textwidth}{@{}>{\bfseries}p{4.3cm} X@{}}
\toprule
Notion & Ce qu'il faut retenir \\
\midrule
Classe & Modèle décrivant des données et des comportements : \py{class Voiture:} \\
Objet / instance & Réalisation concrète de la classe : \py{v = Voiture("Toyota", "Bleu")} \\
Constructeur & \py{__init__}, appelé automatiquement, initialise les attributs \\
\py{self} & Référence à l'objet courant ; premier paramètre de toute méthode d'instance \\
Attribut d'instance & Donnée propre à l'objet : \py{self.marque} \\
Méthode & Fonction définie dans la classe, appelée par \py{objet.methode()} \\
Attribut privé & \py{self.__prix} : inaccessible depuis l'extérieur (\emph{name mangling}) \\
Accesseur (\emph{getter}) & \py{get_prix()} : lit un attribut privé \\
Mutateur (\emph{setter}) & \py{set_prix(v)} : écrit un attribut privé après validation \\
Héritage & \py{class Voiture(Vehicule)} : la fille reçoit les membres de la mère \\
\py{super()} & Accède à la classe mère, notamment \py{super().__init__(...)} \\
Redéfinition & La fille réécrit une méthode de la mère (\py{faire_son}) \\
\bottomrule
\end{tabularx}
\end{center}

\begin{piege}
\textbf{Erreurs les plus fréquentes.}
oublier \py{self} dans la définition d'une méthode ;
confondre \py{return} (retourne une valeur) et \py{print} (affiche) ;
écrire \py{objet.__attribut} depuis l'extérieur ;
oublier \py{super().__init__(...)} dans le constructeur d'une classe fille ;
appeler une méthode sans les parenthèses (\py{v.afficher_details} retourne la méthode,
pas son résultat).
\end{piege}

\end{document}
"
\documentclass[11pt,a4paper]{article}

\usepackage[T1]{fontenc}
\usepackage[utf8]{inputenc}
\usepackage{lmodern}
\usepackage[french]{babel}
\usepackage[a4paper,margin=1.8cm,bottom=2.2cm]{geometry}
\usepackage[table]{xcolor}
\usepackage{booktabs,tabularx,array}
\usepackage{enumitem}
\usepackage{needspace}
\usepackage{fancyhdr}
\usepackage{amssymb,amsmath}
\usepackage[most]{tcolorbox}
\usepackage{listings}
\usepackage[strings]{underscore}  % permet d'écrire __init__ sans échappement
\usepackage{hyperref}
\hypersetup{hidelinks}

\providecommand{\modecorrige}{0}
\newif\ifcorrige
\ifnum\modecorrige=1 \corrigetrue \else \corrigefalse \fi

% Palette alignée sur les autres supports du dépôt
\definecolor{ipred}{HTML}{C0392B}
\definecolor{cdcbar}{HTML}{A67C3D}
\definecolor{cdcbg}{HTML}{F4F1EA}
\definecolor{qbar}{HTML}{4A6B5D}
\definecolor{qbg}{HTML}{EDF1EF}
\definecolor{corr}{HTML}{1F5FA8}
\definecolor{corrbg}{HTML}{E8F0FA}
\definecolor{grisfoot}{HTML}{828282}
\definecolor{codebg}{HTML}{F7F7F4}
\definecolor{codekw}{HTML}{8E44AD}
\definecolor{codestr}{HTML}{1F7A4D}
\definecolor{codecom}{HTML}{8A8A8A}

% ---------------------------------------------------------------- code Python
\lstdefinestyle{py}{
  language=Python,
  basicstyle=\ttfamily\small,
  keywordstyle=\color{codekw}\bfseries,
  stringstyle=\color{codestr},
  commentstyle=\color{codecom}\itshape,
  showstringspaces=false,
  backgroundcolor=\color{codebg},
  frame=single,
  rulecolor=\color{black!25},
  framesep=5pt,
  xleftmargin=4pt,
  xrightmargin=4pt,
  columns=fullflexible,
  keepspaces=true,
  breaklines=true,
  extendedchars=true,
  literate={é}{{\'e}}1 {è}{{\`e}}1 {ê}{{\^e}}1 {à}{{\`a}}1 {ç}{{\c c}}1
           {ù}{{\`u}}1 {û}{{\^u}}1 {î}{{\^\i}}1 {ô}{{\^o}}1 {É}{{\'E}}1 {À}{{\`A}}1 {È}{{\`E}}1 {Ç}{{\c C}}1
           {«}{{\guillemotleft}}1 {»}{{\guillemotright}}1,
  morekeywords={self,True,False,None},
}
\lstset{style=py}
\newcommand{\py}[1]{\texttt{#1}}

% --------------------------------------------------------------- encadrés
\newtcolorbox{methode}[1][]{enhanced,breakable,colback=cdcbg,colframe=cdcbar,boxrule=0pt,
  leftrule=3pt,arc=0pt,outer arc=0pt,left=6pt,right=6pt,top=4pt,bottom=4pt,
  title={#1},coltitle=cdcbar,fonttitle=\bfseries,detach title,
  before upper={\tcbtitle\par\smallskip}}
\newtcolorbox{rappel}[1][]{enhanced,breakable,colback=qbg,colframe=qbar,boxrule=0pt,
  leftrule=3pt,arc=0pt,outer arc=0pt,left=6pt,right=6pt,top=4pt,bottom=4pt,
  title={#1},coltitle=qbar,fonttitle=\bfseries,attach title to upper={\ :\ }}
\newtcolorbox{piege}{enhanced,breakable,colback=red!4,colframe=ipred,boxrule=0pt,
  leftrule=3pt,arc=0pt,outer arc=0pt,left=6pt,right=6pt,top=4pt,bottom=4pt}

% --------------------------------------------------------------- exercices
\newcounter{exo}[section]
\renewcommand{\theexo}{\arabic{exo}}
\newtcolorbox{quest}{enhanced,breakable,colback=qbg,colframe=qbar,boxrule=0pt,
  leftrule=3pt,arc=0pt,outer arc=0pt,left=6pt,right=6pt,top=4pt,bottom=4pt}
\newcommand{\exo}[2][]{% [sous-titre] {énoncé}
  \Needspace{5cm}\stepcounter{exo}%
  \begin{quest}\textbf{Exercice \theexo\ifx\relax#1\relax\else\ — #1\fi.}
  \ifx\relax#2\relax\else\ #2\fi\end{quest}}

% Réponse courte : pointillés dans le sujet, texte bleu dans le corrigé
\newcommand{\rep}[1]{%
  \ifcorrige
    \par\smallskip{\color{corr}\small#1}\par\smallskip
  \else
    \par\nobreak\vspace{1mm}\noindent\color{black!55}\dotfill\par
    \noindent\color{black!55}\dotfill\par\color{black}\vspace{1mm}
  \fi}

% Zone de code vierge (sujet uniquement) de hauteur donnée
\newcommand{\zonecode}[1][4cm]{%
  \ifcorrige\else
    \par\smallskip\noindent
    \begin{tcolorbox}[colback=white,colframe=black!45,boxrule=0.5pt,arc=1pt,height=#1]\end{tcolorbox}%
  \fi}

% Bloc visible seulement dans le corrigé
\newenvironment{solution}
  {\ifcorrige
     \begin{tcolorbox}[enhanced,breakable,colback=corrbg,colframe=corr,boxrule=0.6pt,arc=1pt,
       left=5pt,right=5pt,top=4pt,bottom=4pt,
       title={Corrigé},coltitle=corr,fonttitle=\bfseries\small,detach title,
       before upper={\tcbtitle\par\smallskip},fontupper=\small]%
   \else
     \setbox0\vbox\bgroup
   \fi}
  {\ifcorrige\end{tcolorbox}\else\egroup\fi}

\pagestyle{fancy}
\fancyhf{}
\renewcommand{\headrulewidth}{0pt}
\fancyfoot[L]{\small\itshape\color{grisfoot}TP POO Python\ifcorrige\ — CORRIGÉ\fi}
\fancyfoot[C]{\small\color{grisfoot}BTS CIEL}
\fancyfoot[R]{\small\color{grisfoot}\thepage}

\setlist{itemsep=2pt,topsep=3pt}
\setlength{\parindent}{0pt}
\setlength{\parskip}{4pt}
\renewcommand{\arraystretch}{1.3}

\begin{document}

\begin{center}
{\LARGE\bfseries Programmation orientée objet en Python\par}
\smallskip
{\large Classes, attributs, méthodes --- Encapsulation et héritage\par}
\medskip
BTS CIEL --- Travaux pratiques
\ifcorrige\par\textcolor{corr}{\textbf{Version enseignant --- corrigés intégrés}}\fi
\end{center}

\medskip
\noindent\begin{tabularx}{\textwidth}{@{}>{\bfseries}l X@{}}
Nom, prénom : & \ifcorrige\textcolor{corr}{corrigé}\else\dotfill\fi \\
Classe : & \ifcorrige\textcolor{corr}{---}\else\dotfill\fi \\
Date : & \ifcorrige\textcolor{corr}{---}\else\dotfill\fi \\
\end{tabularx}

\begin{methode}[Objectifs du TP]
\begin{itemize}
  \item Lire une classe Python et identifier son constructeur, ses attributs et ses méthodes.
  \item Instancier un objet et appeler ses méthodes.
  \item Compléter une classe existante en ajoutant une méthode.
  \item Protéger les données d'un objet par l'encapsulation (attributs privés, accesseurs, mutateurs).
  \item Écrire un mutateur qui valide la donnée avant de l'enregistrer.
  \item Construire une classe fille par héritage, redéfinir une méthode et appeler le constructeur parent.
\end{itemize}
\textbf{Consignes} : chaque question se traite d'abord sur le papier, puis se vérifie dans
l'interpréteur Python (fichier \py{tp\_poo.py}). Le code rendu doit s'exécuter sans erreur.
\end{methode}

\begin{rappel}[Vocabulaire de référence]
Une \textbf{classe} est un modèle. Un \textbf{objet} (ou \textbf{instance}) est une réalisation
concrète de ce modèle. Un \textbf{attribut} est une donnée portée par l'objet. Une
\textbf{méthode} est une fonction définie dans la classe et appelée sur un objet. Le paramètre
\py{self} désigne l'objet sur lequel la méthode est appelée.
\end{rappel}

% ============================================================================
\section*{Partie A --- Classes, attributs et méthodes}
\addcontentsline{toc}{section}{Partie A}
\setcounter{exo}{0}

% ---------------------------------------------------------------- A1
\exo[La classe \py{Voiture}]{Lire le code suivant, puis répondre aux questions.}

\begin{lstlisting}
class Voiture:
    def __init__(self, marque, couleur):
        self.marque = marque
        self.couleur = couleur

    def afficher_details(self):
        return f"Voiture {self.marque} de couleur {self.couleur}"
\end{lstlisting}

\textbf{Questions}
\begin{enumerate}
  \item Que fait la méthode \py{__init__} ?
    \rep{C'est le \textbf{constructeur}. Il est appelé automatiquement à la création de l'objet
    et initialise les attributs d'instance \py{self.marque} et \py{self.couleur} à partir des
    valeurs passées en paramètres.}
  \item Combien d'attributs possède la classe \py{Voiture} ?
    \rep{Deux attributs d'instance : \py{marque} et \py{couleur}.}
  \item Que retourne la méthode \py{afficher_details} ?
    \rep{Une chaîne de caractères (\py{str}) construite par une f-string, par exemple
    \py{"Voiture Toyota de couleur Bleu"}. Elle \emph{retourne} la chaîne, elle ne l'affiche pas.}
\end{enumerate}

\textbf{À compléter} : créer une instance de \py{Voiture} de marque \py{"Toyota"} et de couleur
\py{"Bleu"}, appeler \py{afficher_details} et afficher le résultat.
\zonecode[3cm]
\begin{solution}
\begin{lstlisting}
ma_voiture = Voiture("Toyota", "Bleu")
print(ma_voiture.afficher_details())
# Voiture Toyota de couleur Bleu
\end{lstlisting}
\py{print} est nécessaire : sans lui, la chaîne est retournée puis perdue.
\end{solution}

% ---------------------------------------------------------------- A2
\exo[Compléter la classe \py{Chien}]{}

\begin{lstlisting}
class Chien:
    def __init__(self, nom, age):
        self.nom = nom
        self.age = age

    # Méthode aboyer qui retourne "Wouf !"
\end{lstlisting}

\textbf{Questions}
\begin{enumerate}
  \item Quels sont les attributs de la classe \py{Chien} ?
    \rep{\py{nom} et \py{age}, deux attributs d'instance.}
  \item Comment appelle-t-on la méthode \py{__init__} dans le contexte de la POO ?
    \rep{Le \textbf{constructeur} (ou initialiseur). C'est une \emph{méthode spéciale}, reconnaissable
    à son double souligné ; on ne l'appelle jamais directement, elle est déclenchée par
    \py{Chien("Rex", 3)}.}
  \item Ajouter une méthode \py{aboyer} qui retourne \py{"Wouf !"}.
  \item Créer un objet \py{Chien} nommé \py{"Rex"}, âgé de 3 ans, et le faire aboyer.
\end{enumerate}
\zonecode[4.5cm]
\begin{solution}
\begin{lstlisting}
class Chien:
    def __init__(self, nom, age):
        self.nom = nom
        self.age = age

    def aboyer(self):
        return "Wouf !"

rex = Chien("Rex", 3)
print(rex.aboyer())   # Wouf !
\end{lstlisting}
Le paramètre \py{self} est obligatoire dans la définition, même si \py{aboyer} n'utilise aucun
attribut : il reçoit l'objet sur lequel la méthode est appelée.
\end{solution}

% ---------------------------------------------------------------- A3
\exo[La classe \py{Livre}]{}

\begin{lstlisting}
class Livre:
    def __init__(self, titre, auteur):
        self.titre = titre
        self.auteur = auteur

    def description(self):
        return f"{self.titre} par {self.auteur}"
\end{lstlisting}

\textbf{Questions}
\begin{enumerate}
  \item Que fait la méthode \py{description} ?
    \rep{Elle retourne une chaîne qui associe le titre et l'auteur du livre, sous la forme
    \py{"titre par auteur"}.}
  \item Combien de paramètres la méthode \py{__init__} attend-elle, en incluant \py{self} ?
    \rep{Trois : \py{self}, \py{titre} et \py{auteur}. À l'appel, seuls deux arguments sont
    fournis : \py{self} est transmis automatiquement par Python.}
  \item Créer un objet \py{Livre} de titre \py{"Le Petit Prince"} et d'auteur
    \py{"Antoine de Saint-Exupéry"}, puis afficher sa description.
\end{enumerate}
\zonecode[3cm]
\begin{solution}
\begin{lstlisting}
livre = Livre("Le Petit Prince", "Antoine de Saint-Exupery")
print(livre.description())
# Le Petit Prince par Antoine de Saint-Exupery
\end{lstlisting}
\end{solution}

% ---------------------------------------------------------------- A4
\exo[La classe \py{Rectangle}]{}

\begin{lstlisting}
class Rectangle:
    def __init__(self, longueur, largeur):
        self.longueur = longueur
        self.largeur = largeur

    # Méthode calculer_aire qui retourne l'aire du rectangle
\end{lstlisting}

\textbf{Questions}
\begin{enumerate}
  \item Quels sont les attributs de la classe \py{Rectangle} ?
    \rep{\py{longueur} et \py{largeur}.}
  \item Quelle formule mathématique faut-il utiliser pour calculer l'aire ?
    \rep{$\text{aire} = \text{longueur} \times \text{largeur}$.}
  \item Ajouter une méthode \py{calculer_aire} qui retourne l'aire du rectangle.
  \item Créer un rectangle de longueur 5 et de largeur 3, puis afficher son aire.
\end{enumerate}
\zonecode[4.5cm]
\begin{solution}
\begin{lstlisting}
class Rectangle:
    def __init__(self, longueur, largeur):
        self.longueur = longueur
        self.largeur = largeur

    def calculer_aire(self):
        return self.longueur * self.largeur

r = Rectangle(5, 3)
print(r.calculer_aire())   # 15
\end{lstlisting}
Une méthode peut aussi être ajoutée pour le périmètre :\\
\py{return 2 * (self.longueur + self.largeur)}.
\end{solution}

% ---------------------------------------------------------------- A5
\exo[Création complète d'une classe \py{Etudiant}]{On souhaite modéliser un étudiant par son
nom et sa note.}

\begin{enumerate}
  \item Créer une classe \py{Etudiant} possédant un constructeur \py{__init__} prenant en
    paramètres \py{nom} et \py{note}, et deux attributs d'instance \py{self.nom} et \py{self.note}.
  \item Ajouter une méthode \py{a_reussi} qui retourne \py{True} si la note est supérieure ou
    égale à 10, et \py{False} sinon.
  \item Tester la classe en créant un étudiant avec une note de 8.
  \item Ajouter une méthode \py{augmenter_note} qui augmente la note de 2 points.
  \item Créer un étudiant nommé \py{"Alice"} avec une note de 15 et afficher si elle a réussi,
    grâce à la méthode \py{a_reussi}.
\end{enumerate}
\zonecode[7cm]
\begin{solution}
\begin{lstlisting}
class Etudiant:
    def __init__(self, nom, note):
        self.nom = nom
        self.note = note

    def a_reussi(self):
        return self.note >= 10

    def augmenter_note(self):
        self.note = self.note + 2

# 3. test avec une note de 8
e = Etudiant("Bob", 8)
print(e.a_reussi())          # False

# 5. Alice avec une note de 15
alice = Etudiant("Alice", 15)
print(alice.a_reussi())      # True
\end{lstlisting}
\textbf{Points d'attention.}
\py{return self.note >= 10} suffit : la comparaison produit déjà un booléen, un
\py{if ... return True else return False} est inutile.
\py{augmenter_note} \emph{modifie} l'attribut et ne retourne rien : elle agit sur l'état de l'objet.
\end{solution}

% ============================================================================
\clearpage
\section*{Partie B --- Encapsulation et héritage}
\addcontentsline{toc}{section}{Partie B}
\setcounter{exo}{0}

\begin{rappel}[Encapsulation]
Préfixer un attribut par deux soulignés (\py{self.__nom}) déclenche le \emph{name mangling} :
Python renomme l'attribut en \py{_Classe__nom}, ce qui le rend inaccessible par son nom
d'origine depuis l'extérieur. On y accède alors par un \textbf{accesseur} (\emph{getter},
\py{get_nom}) et on le modifie par un \textbf{mutateur} (\emph{setter}, \py{set_nom}), ce qui
permet de contrôler les valeurs acceptées.
\end{rappel}

\begin{rappel}[Héritage]
\py{class Chat(Animal)} signifie que \py{Chat} \textbf{hérite} de \py{Animal} : la classe fille
reçoit les attributs et méthodes de la classe mère. Elle peut en ajouter, ou
\textbf{redéfinir} une méthode existante (\emph{override}). \py{super().__init__(...)} appelle le
constructeur de la classe mère.
\end{rappel}

% ---------------------------------------------------------------- B1
\exo[Encapsulation simple avec attribut privé]{}

\begin{lstlisting}
class Personne:
    def __init__(self, nom, age):
        self.__nom = nom
        self.__age = age

    def get_nom(self):
        return self.__nom

    # Méthode get_age()
\end{lstlisting}

\textbf{Questions}
\begin{enumerate}
  \item Pourquoi les attributs \py{__nom} et \py{__age} sont-ils privés ?
    \rep{Pour empêcher leur modification directe depuis l'extérieur de la classe. L'accès passe
    par des méthodes, ce qui permet de contrôler la lecture et l'écriture et de modifier
    l'implémentation interne sans casser le code qui utilise la classe.}
  \item Comment accède-t-on à l'attribut \py{__nom} ?
    \rep{Par l'accesseur : \py{p.get_nom()}. L'écriture \py{p.__nom} provoque une
    \py{AttributeError}.}
  \item Ajouter une méthode \py{get_age} pour accéder à l'attribut \py{__age}.
  \item Créer une instance de \py{Personne} de nom \py{"Luc"} et d'âge 25, puis afficher son nom.
\end{enumerate}
\zonecode[4.5cm]
\begin{solution}
\begin{lstlisting}
class Personne:
    def __init__(self, nom, age):
        self.__nom = nom
        self.__age = age

    def get_nom(self):
        return self.__nom

    def get_age(self):
        return self.__age

p = Personne("Luc", 25)
print(p.get_nom())    # Luc
print(p.get_age())    # 25
\end{lstlisting}
\end{solution}

% ---------------------------------------------------------------- B2
\exo[Ajouter un mutateur]{}

\begin{lstlisting}
class CompteBancaire:
    def __init__(self, titulaire, solde):
        self.__titulaire = titulaire
        self.__solde = solde

    def get_solde(self):
        return self.__solde
\end{lstlisting}

\textbf{Questions}
\begin{enumerate}
  \item Que signifie l'encapsulation dans cette classe ?
    \rep{Le titulaire et le solde sont privés : aucun code extérieur ne peut écrire
    \py{compte.__solde = 1000000}. Le solde ne peut être consulté que par \py{get_solde} et
    modifié que par les méthodes que la classe met à disposition.}
  \item Que fait la méthode \py{get_solde} ?
    \rep{Elle retourne la valeur de l'attribut privé \py{__solde} : c'est un accesseur en
    lecture seule.}
  \item Créer un compte bancaire pour \py{"Emma"} avec un solde de 500 et afficher le solde.
  \item Ajouter une méthode \py{set_solde} pour modifier le solde.
\end{enumerate}
\zonecode[5cm]
\begin{solution}
\begin{lstlisting}
class CompteBancaire:
    def __init__(self, titulaire, solde):
        self.__titulaire = titulaire
        self.__solde = solde

    def get_solde(self):
        return self.__solde

    def set_solde(self, nouveau_solde):
        self.__solde = nouveau_solde

compte = CompteBancaire("Emma", 500)
print(compte.get_solde())   # 500
compte.set_solde(750)
print(compte.get_solde())   # 750
\end{lstlisting}
\textbf{Prolongement.} Un mutateur brut n'apporte pas de sécurité. Une version plus réaliste
refuse un solde négatif, ou expose plutôt des opérations métier :
\py{deposer(montant)} et \py{retirer(montant)} avec vérification de la provision.
\end{solution}

% ---------------------------------------------------------------- B3
\exo[Encapsulation avec validation]{}

\begin{lstlisting}
class Etudiant:
    def __init__(self, nom, note):
        self.__nom = nom
        self.__note = note

    def get_note(self):
        return self.__note

    # Méthode set_note qui accepte uniquement des notes entre 0 et 20
\end{lstlisting}

\textbf{Questions}
\begin{enumerate}
  \item Pourquoi utiliser un mutateur pour la note ?
    \rep{Pour garantir un \emph{invariant} : une note reste comprise entre 0 et 20. Le mutateur
    est le seul point d'entrée en écriture, il peut donc vérifier la valeur avant de
    l'enregistrer et refuser les valeurs incohérentes.}
  \item Que se passe-t-il si on essaie d'accéder directement à \py{__note} ?
    \rep{Python lève une \py{AttributeError} : à cause du \emph{name mangling}, l'attribut
    s'appelle en réalité \py{_Etudiant__note} et le nom \py{__note} n'existe pas depuis
    l'extérieur.}
  \item Ajouter une méthode \py{set_note} qui n'accepte que des notes entre 0 et 20.
  \item Créer un étudiant nommé \py{"Paul"} avec une note de 15, puis essayer de changer la note
    à 25.
\end{enumerate}
\zonecode[5.5cm]
\begin{solution}
\begin{lstlisting}
class Etudiant:
    def __init__(self, nom, note):
        self.__nom = nom
        self.__note = note

    def get_note(self):
        return self.__note

    def set_note(self, nouvelle_note):
        if 0 <= nouvelle_note <= 20:
            self.__note = nouvelle_note
        else:
            print("Note invalide : elle doit etre comprise entre 0 et 20.")

paul = Etudiant("Paul", 15)
print(paul.get_note())   # 15
paul.set_note(25)        # Note invalide : ...
print(paul.get_note())   # 15 : la note n'a pas change
\end{lstlisting}
\textbf{Variante professionnelle.} Au lieu d'un \py{print}, lever une exception rend l'erreur
impossible à ignorer : \py{raise ValueError("note hors intervalle")}.
\end{solution}

% ---------------------------------------------------------------- B4
\exo[Héritage simple]{}

\begin{lstlisting}
class Animal:
    def __init__(self, nom):
        self.__nom = nom

    def get_nom(self):
        return self.__nom

    def faire_son(self):
        return "Son generique"

class Chat(Animal):
    # Méthode faire_son pour retourner "Miaou"
\end{lstlisting}

\textbf{Questions}
\begin{enumerate}
  \item Quelle classe est la classe parent ?
    \rep{\py{Animal}. \py{Chat} est la classe fille (ou classe dérivée).}
  \item Que fait l'héritage dans ce cas ?
    \rep{\py{Chat} récupère sans les réécrire le constructeur \py{__init__}, l'attribut privé
    \py{__nom} et la méthode \py{get_nom} d'\py{Animal}. Elle peut en plus redéfinir
    \py{faire_son} pour lui donner son comportement propre.}
  \item Redéfinir la méthode \py{faire_son} pour qu'elle retourne \py{"Miaou"}.
  \item Créer un \py{Chat} nommé \py{"Minou"} et afficher son son.
\end{enumerate}
\zonecode[4.5cm]
\begin{solution}
\begin{lstlisting}
class Chat(Animal):
    def faire_son(self):
        return "Miaou"

minou = Chat("Minou")
print(minou.get_nom())     # Minou  (methode heritee)
print(minou.faire_son())   # Miaou  (methode redefinie)
\end{lstlisting}
\py{Chat} ne définit pas de \py{__init__} : celui d'\py{Animal} est utilisé tel quel.
Appeler \py{faire_son()} sur un \py{Chat} exécute la version de \py{Chat} : c'est le
\textbf{polymorphisme}.
\end{solution}

% ---------------------------------------------------------------- B5
\exo[Héritage avec nouvel attribut]{}

\begin{lstlisting}
class Vehicule:
    def __init__(self, marque):
        self.__marque = marque

    def get_marque(self):
        return self.__marque

class Voiture(Vehicule):
    def __init__(self, marque, couleur):
        super().__init__(marque)
        self.__couleur = couleur

    # À compléter : ajoutez une méthode get_couleur
\end{lstlisting}

\textbf{Questions}
\begin{enumerate}
  \item Que fait \py{super().__init__(marque)} ?
    \rep{Il appelle le constructeur de la classe mère \py{Vehicule} pour initialiser la partie
    héritée de l'objet, ici l'attribut privé \py{__marque}. Sans cet appel, \py{get_marque}
    échouerait car l'attribut n'aurait jamais été créé.}
  \item Combien d'attributs privés possède la classe \py{Voiture} ?
    \rep{Deux au total dans l'objet : \py{_Vehicule__marque}, hérité et initialisé par
    \py{super()}, et \py{_Voiture__couleur}, propre à \py{Voiture}. La classe \py{Voiture} n'en
    déclare qu'un seul elle-même.}
  \item À compléter : ajouter une méthode \py{get_couleur}.
\end{enumerate}
\zonecode[4cm]
\begin{solution}
\begin{lstlisting}
class Voiture(Vehicule):
    def __init__(self, marque, couleur):
        super().__init__(marque)
        self.__couleur = couleur

    def get_couleur(self):
        return self.__couleur

v = Voiture("Renault", "Bleu")
print(v.get_marque())    # Renault
print(v.get_couleur())   # Bleu
\end{lstlisting}
\begin{piege}
\py{__marque} est privé \emph{à \py{Vehicule}} : à cause du \emph{name mangling},
une méthode de \py{Voiture} ne peut pas écrire \py{self.__marque}. Elle doit passer par
\py{self.get_marque()}. C'est exactement le rôle de l'accesseur hérité.
\end{piege}
\end{solution}

% ---------------------------------------------------------------- B6
\exo[Encapsulation et mutateur avec validation]{}

\begin{lstlisting}
class Livre:
    def __init__(self, titre, prix):
        self.__titre = titre
        self.__prix = prix

    def get_prix(self):
        return self.__prix

    # À compléter : ajoutez une méthode set_prix qui refuse les prix négatifs
\end{lstlisting}

\textbf{Questions}
\begin{enumerate}
  \item Pourquoi le prix est-il un attribut privé ?
    \rep{Parce qu'un prix doit respecter une règle métier : il ne peut pas être négatif. Le
    rendre privé interdit l'écriture directe et oblige à passer par une méthode qui applique
    cette règle.}
  \item Que fait un mutateur avec validation ?
    \rep{Il vérifie la valeur reçue avant de l'affecter à l'attribut. Si la valeur est valide,
    il met l'attribut à jour ; sinon il refuse la modification (message d'erreur ou exception)
    et l'objet conserve son état précédent, qui reste cohérent.}
  \item Créer un \py{Livre} intitulé \py{"Python 101"} avec un prix de 30, puis essayer de
    définir un prix de $-10$.
\end{enumerate}
\zonecode[5.5cm]
\begin{solution}
\begin{lstlisting}
class Livre:
    def __init__(self, titre, prix):
        self.__titre = titre
        self.__prix = prix

    def get_prix(self):
        return self.__prix

    def set_prix(self, nouveau_prix):
        if nouveau_prix >= 0:
            self.__prix = nouveau_prix
        else:
            print("Prix invalide : un prix ne peut pas etre negatif.")

livre = Livre("Python 101", 30)
print(livre.get_prix())   # 30
livre.set_prix(-10)       # Prix invalide : ...
print(livre.get_prix())   # 30 : le prix est inchange
\end{lstlisting}
\textbf{Remarque.} Le constructeur devrait appliquer la même validation, par exemple en
appelant \py{self.set_prix(prix)} au lieu d'affecter directement : sinon
\py{Livre("X", -10)} crée un objet invalide dès le départ.
\end{solution}

% ============================================================================
\clearpage
\section*{Synthèse}

\begin{center}
\begin{tabularx}{\textwidth}{@{}>{\bfseries}p{4.3cm} X@{}}
\toprule
Notion & Ce qu'il faut retenir \\
\midrule
Classe & Modèle décrivant des données et des comportements : \py{class Voiture:} \\
Objet / instance & Réalisation concrète de la classe : \py{v = Voiture("Toyota", "Bleu")} \\
Constructeur & \py{__init__}, appelé automatiquement, initialise les attributs \\
\py{self} & Référence à l'objet courant ; premier paramètre de toute méthode d'instance \\
Attribut d'instance & Donnée propre à l'objet : \py{self.marque} \\
Méthode & Fonction définie dans la classe, appelée par \py{objet.methode()} \\
Attribut privé & \py{self.__prix} : inaccessible depuis l'extérieur (\emph{name mangling}) \\
Accesseur (\emph{getter}) & \py{get_prix()} : lit un attribut privé \\
Mutateur (\emph{setter}) & \py{set_prix(v)} : écrit un attribut privé après validation \\
Héritage & \py{class Voiture(Vehicule)} : la fille reçoit les membres de la mère \\
\py{super()} & Accède à la classe mère, notamment \py{super().__init__(...)} \\
Redéfinition & La fille réécrit une méthode de la mère (\py{faire_son}) \\
\bottomrule
\end{tabularx}
\end{center}

\begin{piege}
\textbf{Erreurs les plus fréquentes.}
oublier \py{self} dans la définition d'une méthode ;
confondre \py{return} (retourne une valeur) et \py{print} (affiche) ;
écrire \py{objet.__attribut} depuis l'extérieur ;
oublier \py{super().__init__(...)} dans le constructeur d'une classe fille ;
appeler une méthode sans les parenthèses (\py{v.afficher_details} retourne la méthode,
pas son résultat).
\end{piege}

\end{document}
